\motto{Algebra is the offer made by the devil to the mathematician. The devil says: ``I will give you this powerful machine, it will answer any question you like. All you need to do is give me your soul: give up geometry and you will have this marvelous machine.''

--- Michael Atiyah\footnote{Sir Michael Francis Atiyah (1929--2019) wrote the influential \emph{Introduction to commutative algebra} together with I.~G.~MacDonald, whose name is often omitted or misspelled.}}
\chapter{Comparison of singularities}
\label{chap:comp}

A central theme of this book is that singularities of quasi-plurisubharmonic functions can be compared at several levels of precision. The preorder $\preceq$ of \cref{def:parorder} compares potentials up to additive constants. It is the most concrete notion, but it is often too fine for global pluripotential theory: the $P$-envelope identifies potentials with the same model representative, while multiplier ideals identify singularities that have the same integrability behavior. This chapter makes these two coarser comparisons systematic.

We first introduce the $P$-preorder and the $\mathcal I$-preorder on $\QPSH(X)$. The $P$-preorder is defined by comparing $P$-envelopes, and its equivalence classes are the $P$-singularity types. The $\mathcal I$-preorder is defined by comparing multiplier ideals $\mathcal I(k\varphi)$ for all $k\in\mathbb Z_{>0}$. These two preorders are compatible with ordinary singularity comparison, but they capture different information: the former is governed by non-pluripolar mass, while the latter is governed by valuative and multiplier-ideal data.

The second main object of the chapter is the pseudometric $d_S$. For positive-mass potentials in a fixed class, $d_S$ is defined through the Darvas $d_1$-distance between the associated geodesic rays. Its zero locus is precisely $P$-equivalence. Thus $d_S$ should be viewed as a convergence structure on $P$-singularity types. The chapter proves the basic properties of $d_S$, including monotone convergence, sandwich criteria and compatibility with changing the background form.

The final results establish the continuity properties that will be used throughout the rest of the book. We prove continuity of the $\mathcal I$-envelope, continuity of non-pluripolar mass and volume, and a general equisingularity criterion. These results allow one to pass from analytic approximations to limiting singularities while controlling both $P$-types and multiplier ideals. They are used repeatedly in the theory of $\mathcal I$-good singularities, trace operators, test curves, partial Okounkov bodies and non-Archimedean potentials.

\section{The \texorpdfstring{$P$}{P} and \texorpdfstring{$\mathcal{I}$}{I}-preorders}\label{sec:PIpartialorder}
Let $X$ be a connected compact Kähler manifold of dimension $n$.


Recall that we defined a preorder on $\QPSH(X)$ in \cref{def:parorder} to compare the singularity types of quasi-plurisubharmonic functions. The problem with this preorder is that it is too fine. For our purposes, it is useful to consider rougher relations.

\subsection{The definitions of the preorders}
\label{subsec:preorder_def}

Recall that the $P$-envelope is defined in \cref{def:Penv}.
\begin{definition}\label{def:Pmoresing}
    For $\varphi,\psi\in \QPSH(X)$, we say $\varphi$ is \emph{$P$-more singular than $\psi$}\index{P-more singular@$P$-more singular} and write $\varphi\preceq_P \psi$\index{$\varphi\preceq_P \psi$} if for some (equivalently, every, due to \cref{lma:Pproj_insens_omega} below) closed smooth real $(1,1)$-form $\theta$ on $X$ such that $\varphi,\psi\in \PSH(X,\theta)_{>0}$, we have
        \begin{equation}\label{eq:Penvleq}
        P_{\theta}[\varphi]\leq P_{\theta}[\psi].
        \end{equation}
    If $\varphi\preceq_P \psi$ and $\psi \preceq_P \varphi$, we write $\varphi\sim_P \psi$ and say that $\varphi$ and $\psi$ have the same \emph{$P$-singularity type}\index{P-singularity type@$P$-singularity type}.
\end{definition}
Recall that $\PSH(X,\theta)_{>0}$ was introduced in \eqref{eq:PSHpos}.
Note that if $\varphi\preceq \psi$, then $\varphi\preceq_P \psi$. So the $P$-preorder is \emph{coarser} than $\preceq$.

The $P$-equivalence relation can also be formulated using the language of relative full mass classes developed in \cref{subsec:fullmass}.
In fact, $\varphi,\psi\in \PSH(X,\theta)_{>0}$ satisfies $\varphi\sim_P \psi$ if and only if $\varphi,\psi \in \mathcal{E}(X,\theta;\phi)$ for a common model potential $\phi\in \PSH(X,\theta)_{>0}$.


The condition \eqref{eq:Penvleq} is independent of the choice of $\theta$:
\begin{lemma}\label{lma:Pproj_insens_omega}
    Let $\varphi,\psi\in \PSH(X,\theta)_{>0}$. For any Kähler form $\omega$ on $X$, the following are equivalent:
    \begin{enumerate}
        \item\label{itm:lma:Pproj_insens_omega:1} $P_{\theta}[\varphi]\leq P_{\theta}[\psi]$;
        \item\label{itm:lma:Pproj_insens_omega:2} $P_{\theta+\omega}[\varphi]\leq P_{\theta+\omega}[\psi]$.
    \end{enumerate}
\end{lemma}
In particular, $\preceq_P$ defines a preorder on $\QPSH(X)$.
\begin{proof}
    \ref{itm:lma:Pproj_insens_omega:1} $\implies$ \ref{itm:lma:Pproj_insens_omega:2}. Assume \ref{itm:lma:Pproj_insens_omega:1}. Observe that
    \[
    \varphi\preceq P_{\theta}[\varphi]\leq P_{\theta+\omega}[\varphi].
    \]
    Taking the $P_{\theta+\omega}$-envelope, we find
    \[
    P_{\theta+\omega}[\varphi] \leq P_{\theta+\omega}\Bigl[P_{\theta}[\varphi]\Bigr] \leq P_{\theta+\omega}\Bigl[P_{\theta+\omega}[\varphi]\Bigr]=P_{\theta+\omega}[\varphi],
    \]
    where the equality follows from \cref{thm:Pvarphidiffdef}. In particular,
    \begin{equation}\label{eq:doubleP}
    P_{\theta+\omega}[\varphi]=P_{\theta+\omega}\Bigl[P_{\theta}[\varphi]\Bigr].
    \end{equation}
    A similar formula holds for $\psi$. So we see that \ref{itm:lma:Pproj_insens_omega:2} holds.

    \ref{itm:lma:Pproj_insens_omega:2} $\implies$ \ref{itm:lma:Pproj_insens_omega:1}. Assume \ref{itm:lma:Pproj_insens_omega:2}. By \eqref{eq:doubleP}, we may assume that $\varphi$ and $\psi$ are both model potentials in $\PSH(X,\theta)_{>0}$.

    Observe that $\varphi\lor \psi\preceq P_{\theta+\omega}[\psi]$.
    It follows that $P_{\theta+\omega}[\varphi\lor \psi]\leq P_{\theta+\omega}[\psi]$.
    The reverse inequality is trivial, so
    \[
    P_{\theta+\omega}[\varphi\lor \psi]= P_{\theta+\omega}[\psi].
    \]
    From the direction we have proved, for any $C\geq 1$,
    \[
    P_{\theta+C\omega}[\varphi\lor \psi]= P_{\theta+C\omega}[\psi].
    \]
    So by \cref{prop:Ppresmass},
    \[
    \int_X \Bigl(\theta+C\omega+\ddc(\varphi\lor \psi)\Bigr)^n=\int_X\left(\theta+C\omega+\ddc\psi\right)^n.
    \]
    Since both sides are polynomials in $C$, the equality extends to $C=0$, namely,
    \[
    \int_X \theta_{\varphi\lor \psi}^n=\int_X \theta_{\psi}^n.
    \]
    In particular, $\varphi\lor \psi\leq P_{\theta}[\psi]=\psi$ by \eqref{eq:Penvdef}. So \ref{itm:lma:Pproj_insens_omega:1} follows.
\end{proof}
As a consequence of \cref{lma:Pproj_insens_omega}, we can define the $P$-preorder at the level of currents.
Given closed positive $(1,1)$-currents $T=\theta_{\varphi}$, $S=\theta'_{\psi}$, we write $T\preceq_P S$ (resp. $T\sim_P S$) if $\varphi\preceq_P \psi$ (resp. $\varphi\sim_P \psi$). This definition is independent of the decompositions of $T$ and $S$.

As a first example of $P$-equivalence, we have:
\begin{example}\label{ex:Pequiv}
    Let $\theta$ be a closed smooth real $(1,1)$-form on $X$ and $\varphi\in \PSH(X,\theta)_{>0}$. Then
    \[
        \varphi\sim_P P_{\theta}[\varphi]\footnotemark.
    \]
    \footnotetext{It is not known to us whether the same holds when $\varphi$ has vanishing mass.}
    This follows immediately from \cref{thm:Pvarphidiffdef}.
\end{example}

We give a very useful criterion of the $P$-equivalence in terms of the non-pluripolar masses.
\begin{proposition}\label{prop:Pequivchar2}
    Let $\theta$ be a closed smooth real $(1,1)$-form on $X$ representing a big cohomology class.
    Let $\varphi,\psi\in \PSH(X,\theta)$ and $\varphi\preceq_P \psi$. Then the following are equivalent:
    \begin{enumerate}
        \item\label{itm:prop:Pequivchar2:1} $\varphi\sim_P \psi$;
        \item\label{itm:prop:Pequivchar2:2} for each $j=0,\ldots,n$, we have
            \begin{equation}\label{eq:mixedmassequal}
                \int_X \theta_{\varphi}^j\wedge \theta_{V_{\theta}}^{n-j}=\int_X \theta_{\psi}^j\wedge \theta_{V_{\theta}}^{n-j}.
            \end{equation}
    \end{enumerate}
    Assume furthermore that $\varphi,\psi\in \PSH(X,\theta)_{>0}$, then these conditions are equivalent to the following:
    \begin{enumerate}[resume]
        \item\label{itm:prop:Pequivchar2:3} We have
            \[
                \int_X \theta_{\varphi}^n=\int_X \theta_{\psi}^n.
            \]
    \end{enumerate}
\end{proposition}
Recall that $V_{\theta}$ is introduced in \eqref{eq:Vtheta}.

\begin{proof}
    We first prove the equivalence between \ref{itm:prop:Pequivchar2:1} and \ref{itm:prop:Pequivchar2:3} when $\varphi,\psi\in \PSH(X,\theta)_{>0}$.

    \ref{itm:prop:Pequivchar2:1} $\implies$ \ref{itm:prop:Pequivchar2:3}. Assume that $\varphi\sim_P \psi$. By \cref{lma:Pproj_insens_omega}, we have
    \[
        P_{\theta}[\varphi]=P_{\theta}[\psi].
    \]
    So \ref{itm:prop:Pequivchar2:3} follows from \cref{prop:Ppresmass}.

    \ref{itm:prop:Pequivchar2:3} $\implies$ \ref{itm:prop:Pequivchar2:1}. Since $\varphi\preceq_P\psi$, we have $P_{\theta}[\varphi]\leq P_{\theta}[\psi]$.
    Moreover, \cref{prop:Ppresmass} shows that $P_{\theta}[\varphi]$ and $P_{\theta}[\psi]$ have the same positive mass.
    By \cref{thm:Pvarphidiffdef}, $P_{\theta}[\psi]$ is a candidate in the characterization of $P_{\theta}\bigl[P_{\theta}[\varphi]\bigr]=P_{\theta}[\varphi]$.
    Hence $P_{\theta}[\psi]\leq P_{\theta}[\varphi]$, and therefore $P_{\theta}[\varphi]=P_{\theta}[\psi]$.
    So \ref{itm:prop:Pequivchar2:1} follows.

    Let us come back to the general case with $\varphi,\psi\in \PSH(X,\theta)$.

    \ref{itm:prop:Pequivchar2:1} $\implies$ \ref{itm:prop:Pequivchar2:2}.
    Fix $j\in \{0,\ldots,n\}$. We argue \eqref{eq:mixedmassequal}.

    Take a Kähler form $\omega$ on $X$. By \cref{lma:Pproj_insens_omega}, for each $\epsilon>0$, we have
    \[
        P_{\theta+\epsilon\omega}[\varphi]=  P_{\theta+\epsilon\omega}[\psi].
    \]
    It follows from \cref{prop:Ppresmass} that
    \[
        \begin{aligned}
        \int_X \left(\theta+\epsilon\omega+\ddc \psi  \right)^j\wedge \theta_{V_{\theta}}^{n-j}=&\int_X \Bigl(\theta+\epsilon\omega+\ddc P_{\theta+\epsilon\omega}[\psi]  \Bigr)^j\wedge \theta_{V_{\theta}}^{n-j}\\
        =&\int_X \Bigl(\theta+\epsilon\omega+\ddc P_{\theta+\epsilon\omega}[\varphi]  \Bigr)^j\wedge \theta_{V_{\theta}}^{n-j} \\
        =& \int_X \left(\theta+\epsilon\omega+\ddc \varphi  \right)^j\wedge \theta_{V_{\theta}}^{n-j}.
        \end{aligned}
    \]
    Since the two extremes are both polynomials in $\epsilon$ by the multi-linearity of the non-pluripolar products \itemref{itm:prop:npp1:6}, we conclude that the same holds when $\epsilon=0$, that is, \eqref{eq:mixedmassequal} holds.

    \ref{itm:prop:Pequivchar2:2} $\implies$ \ref{itm:prop:Pequivchar2:1}. Assume \eqref{eq:mixedmassequal} holds for all $j=0,\ldots,n$. Then the same holds with $P_{\theta}[\psi]$ in place of $\psi$, thanks to \cref{prop:Ppresmass}.
    Therefore, for each $t\in (0,1)$, we have
    \[
        \int_X \theta_{t\varphi+(1-t)V_{\theta}}^n=  \int_X \theta_{tP_{\theta}[\psi]+(1-t)V_{\theta}}^n
    \]
    by the binomial expansion. Furthermore, $t\varphi+(1-t)V_{\theta}\preceq tP_{\theta}[\psi]+(1-t)V_{\theta}$.
    By the implication \ref{itm:prop:Pequivchar2:3} $\implies$ \ref{itm:prop:Pequivchar2:1}, we have
    \[
        t\varphi+(1-t)V_{\theta}\sim_P t P_{\theta}[\psi]+(1-t)V_{\theta}
    \]
    for each $t\in (0,1)$.

    Fix a Kähler form $\omega$ on $X$.
    From the implication \ref{itm:prop:Pequivchar2:1} $\implies$ \ref{itm:prop:Pequivchar2:3}, we have
    \[
        \int_X (\theta+\omega)_{t\varphi+(1-t)V_{\theta}}^n=  \int_X (\theta+\omega)_{tP_{\theta}[\psi]+(1-t)V_{\theta}}^n.
    \]
    Since both sides are polynomials in $t$ by the multi-linearity of the non-pluripolar products \itemref{itm:prop:npp1:6}, the same holds when $t=1$. From the implication \ref{itm:prop:Pequivchar2:3} $\implies$ \ref{itm:prop:Pequivchar2:1} again, we have $\varphi\sim_P P_{\theta}[\psi]\sim_P \psi$.
\end{proof}



Next we introduce a different preorder.
\begin{proposition}\label{prop:Iequivchar2}
    Given $\varphi,\psi\in \QPSH(X)$, the following are equivalent:
    \begin{enumerate}
        \item\label{itm:prop:Iequivchar2:1} For any $k\in \mathbb{Z}_{>0}$, we have
        \[
            \mathcal{I}(k\varphi)\subseteq\mathcal{I}(k\psi);
        \]
        \item\label{itm:prop:Iequivchar2:2} for any $\lambda\in \mathbb{R}_{>0}$, we have
            \[
            \mathcal{I}(\lambda\varphi)\subseteq\mathcal{I}(\lambda\psi);
            \]
        \item\label{itm:prop:Iequivchar2:3} for any modification $\pi\colon Y\rightarrow X$ and any $y\in Y$, we have
            \[
                \nu(\pi^*\varphi,y)\geq \nu(\pi^*\psi,y);
            \]
        \item\label{itm:prop:Iequivchar2:4} for any proper bimeromorphic morphism $\pi\colon Y\rightarrow X$ from a Kähler manifold and any $y\in Y$, we have
            \[
                \nu(\pi^*\varphi,y)\geq \nu(\pi^*\psi,y);
            \]
        \item\label{itm:prop:Iequivchar2:5} for any prime divisor $E$ over $X$, we have
            \[
                \nu(\varphi,E)\geq \nu(\psi,E).
            \]
    \end{enumerate}
\end{proposition}
\begin{proof}
    The proof is almost identical to that of \cref{prop:Iequivchar}.
\end{proof}

\begin{definition}\label{def:Imoresing}
    Let $\varphi,\psi\in \QPSH(X)$, we say $\varphi$ is \emph{$\mathcal{I}$-more singular than $\psi$}\index{I-more singular@$\mathcal{I}$-more singular} and write $\varphi\preceq_{\mathcal{I}} \psi$\index{$\varphi\preceq_{\mathcal{I}} \psi$} if the equivalent conditions in \cref{prop:Iequivchar2} are satisfied.
\end{definition}
It is clear that $\preceq_{\mathcal{I}}$ is a preorder on $\QPSH(X)$.

Note that $\varphi\preceq_{\mathcal{I}}\psi$ and $\psi\preceq_{\mathcal{I}}\varphi$ both hold if and only if $\varphi\sim_{\mathcal{I}}\psi$ in the sense of \cref{def:Iequiv}.

Given closed positive $(1,1)$-currents $T=\theta_{\varphi}$, $S=\theta'_{\psi}$, we write $T\preceq_{\mathcal{I}} S$ (resp. $T\sim_{\mathcal{I}} S$) if $\varphi\preceq_{\mathcal{I}} \psi$ (resp. $\varphi\sim_{\mathcal{I}} \psi$). This definition is independent of the decompositions of $T$ and $S$.

\begin{lemma}\label{lma:Plor1}
    Let $\theta$ be a closed smooth real $(1,1)$-form on $X$ representing a big cohomology class.
    Let $\varphi,\psi\in \PSH(X,\theta)_{>0}$. Then
    \begin{equation}\label{eq:Plor1}
    P_{\theta}[\varphi\lor \psi]=P_{\theta}\Bigl[ P_{\theta}[\varphi]\lor P_{\theta}[\psi] \Bigr].
    \end{equation}
\end{lemma}
\begin{proof}
    Since $\varphi\lor \psi\preceq P_{\theta}[\varphi]\lor P_{\theta}[\psi]$, the $\leq$ direction of \eqref{eq:Plor1} follows. Conversely, it suffices to show that
    \[
    P_{\theta}[\varphi\lor \psi]\geq P_{\theta}[\varphi]\lor P_{\theta}[\psi],
    \]
    which follows from $\varphi\preceq\varphi\lor\psi$, $\psi\preceq\varphi\lor\psi$, and the monotonicity of the $P$-envelope with respect to $\preceq$.
\end{proof}

\begin{lemma}\label{lma:reform_preceqP}
    Let $\varphi,\psi\in \QPSH(X)$. Then the following are equivalent:
    \begin{enumerate}
        \item\label{itm:lma:reform_preceqP:1} $\varphi\preceq_P \psi$ (resp. $\varphi\preceq_{\mathcal{I}} \psi$);
        \item\label{itm:lma:reform_preceqP:2} $\varphi\lor \psi\sim_P\psi$ (resp. $\varphi\lor \psi\sim_{\mathcal{I}}\psi$).
    \end{enumerate}
\end{lemma}
\begin{proof}
Take a closed real smooth $(1,1)$-form $\theta$ on $X$ such that $\varphi,\psi\in \PSH(X,\theta)_{>0}$. We only prove the $P$ case, the $\mathcal{I}$ case is similar.

\ref{itm:lma:reform_preceqP:2} $\implies$ \ref{itm:lma:reform_preceqP:1}. Assume \ref{itm:lma:reform_preceqP:2}. By \ref{itm:lma:reform_preceqP:2} and \cref{ex:Pequiv}, $P_{\theta}[\varphi\lor \psi]=P_{\theta}[\psi]\sim_P \psi$. But $\varphi\preceq P_{\theta}[\varphi\lor \psi]$, so \ref{itm:lma:reform_preceqP:1} follows.

\ref{itm:lma:reform_preceqP:1} $\implies$ \ref{itm:lma:reform_preceqP:2}. Assume \ref{itm:lma:reform_preceqP:1}. We may assume that $\varphi,\psi$ are both model by \cref{lma:Plor1}.
    Then $\varphi\leq \psi$ and \ref{itm:lma:reform_preceqP:2} follows.
\end{proof}

\begin{corollary}\label{cor:PimpliesI}
    Let $\varphi,\psi \in \QPSH(X)$. Assume that $\varphi\preceq_P \psi$. Then $\varphi\preceq_{\mathcal{I}}\psi$.
\end{corollary}
Therefore, the preorder $\preceq_{\mathcal{I}}$ is coarser than $\preceq_P$. Given $\varphi,\psi\in \QPSH(X)$, we therefore have the implications:
\[
\varphi\preceq \psi \implies \varphi\preceq_P \psi \implies \varphi\preceq_{\mathcal{I}} \psi.
\]
\begin{proof}
    Thanks to \cref{lma:reform_preceqP}, we have $\varphi\lor \psi\sim_P \psi$. Take a closed smooth real $(1,1)$-form $\theta$ on $X$ such that $\varphi,\psi\in \PSH(X,\theta)_{>0}$.
    Then by \cref{prop:PandPI}, we have
    \[
        \psi\sim_{\mathcal{I}} P_{\theta}[\psi]= P_{\theta}\left[\varphi\lor \psi\right]\sim_{\mathcal{I}} \varphi\lor \psi.
    \]
    Our assertion follows from the $\mathcal{I}$-version of \cref{lma:reform_preceqP}.
\end{proof}

Next we give a few extra characterizations of the $P$-envelope.
\begin{corollary}\label{cor:Pvarphidef3}
    Let $\theta$ be a closed smooth real $(1,1)$-form on $X$ representing a big cohomology class.
    Assume that $\varphi\in \PSH(X,\theta)_{>0}$. Then
    \[
        \begin{aligned}
        P_{\theta}[\varphi]=&\sup\left\{ \psi\in \PSH(X,\theta): \psi\leq 0, \psi\sim_P \varphi \right\}\\
        =&  \sup\left\{ \psi\in \PSH(X,\theta): \psi\leq 0, \psi\preceq_P \varphi \right\}.
        \end{aligned}
    \]
\end{corollary}
Just for comparison, let us recall a few other characterizations of the $P$-envelope for $\varphi\in \PSH(X,\theta)_{>0}$:
\[
\begin{aligned}
    P_{\theta}[\varphi]=&\uscsup\left\{ \psi\in \PSH(X,\theta):\psi\leq 0,\psi\preceq \varphi \right\}\\
    =&\uscsup\left\{ \psi\in \PSH(X,\theta):\psi\leq 0,\psi\sim \varphi \right\}\\
    =&\uscsup[C\in \mathbb{Z}_{>0}] (\varphi+C)\land V_{\theta}\\
    =& \sup\left\{\psi\in \PSH(X,\theta): \psi\leq 0,\varphi\preceq \psi, \int_X \theta_{\varphi}^n=\int_X \theta_{\psi}^n\right\}.
\end{aligned}
\]
\begin{proof}
    Note that $\psi\sim_P \varphi$ implies that $\psi\in \PSH(X,\theta)_{>0}$ by \cref{prop:Ppresmass}. We observe that
    \[
    \begin{split}
    &\sup\left\{ \psi\in \PSH(X,\theta): \psi\leq 0, \psi\sim_P \varphi \right\}\\=&\sup\left\{ \psi\in \PSH(X,\theta): \psi\leq 0, \varphi\preceq \psi, \psi\sim_P \varphi \right\}
    \end{split}
    \]
    since for any candidate $\psi$ for the first supremum, $\psi \lor (\varphi-\sup_X \varphi)$ is larger than $\psi$, and is a candidate for the second supremum thanks to \cref{lma:reform_preceqP}.
    So the first equality is a direct consequence of \cref{prop:Pequivchar2} and \cref{thm:Pvarphidiffdef}.

    Next we prove the second equality. We only need to show that for any $\psi\in \PSH(X,\theta)$ with $\psi\leq 0$ and $\psi\preceq_P \varphi$, we have $\psi\leq P_{\theta}[\varphi]$.

    By \cref{lma:reform_preceqP} and \cref{ex:Pequiv}, we know that $P_{\theta}[\varphi]\lor \psi\sim_P \varphi$ and $P_{\theta}[\varphi]\lor \psi\leq 0$. It follows from the first equality that $\psi\leq P_{\theta}[\varphi]$.
\end{proof}
Similarly, we have a new characterization of the $\mathcal{I}$-envelope.
\begin{corollary}\label{cor:Ienvelopedef2}
    Let $\theta$ be a closed smooth real $(1,1)$-form on $X$ representing a pseudo-effective cohomology class.
    Assume that $\varphi\in \PSH(X,\theta)$. Then
    \[
        P_{\theta}[\varphi]_{\mathcal{I}}=\sup\left\{ \psi\in \PSH(X,\theta): \psi\leq 0, \psi\preceq_{\mathcal{I}} \varphi \right\}.
    \]
\end{corollary}
\begin{proof}
    It suffices to show that for any $\psi\in \PSH(X,\theta)$ with $\psi\leq 0$ and $\psi\preceq_{\mathcal{I}}\varphi$, we have $\psi\leq P_{\theta}[\varphi]_{\mathcal{I}}$.
    By \cref{lma:reform_preceqP} and \cref{prop:Ienvprojection}, we know that $P_{\theta}[\varphi]_{\mathcal{I}}\lor \psi\sim_{\mathcal{I}}\varphi$. Therefore, $P_{\theta}[\varphi]_{\mathcal{I}}\lor \psi$ is a candidate in the definition of $P_{\theta}[\varphi]_{\mathcal{I}}$ (cf. \eqref{eq:Ienvelopedef}), so
    \[
    \psi\leq P_{\theta}[\varphi]_{\mathcal{I}}\lor \psi\leq P_{\theta}[\varphi]_{\mathcal{I}}.
    \]
    We conclude the proof.
\end{proof}

\begin{proposition}\label{prop:Icomparandenvelope}
    Let $\theta$ be a closed real smooth $(1,1)$-form on $X$ and $\varphi,\psi\in \PSH(X,\theta)$. Then the following are equivalent:
    \begin{enumerate}
        \item\label{itm:prop:Icomparandenvelope:1} $\varphi\preceq_{\mathcal{I}}\psi$;
        \item\label{itm:prop:Icomparandenvelope:2} $P_{\theta}[\varphi]_{\mathcal{I}}\leq P_{\theta}[\psi]_{\mathcal{I}}$.
    \end{enumerate}
\end{proposition}
\begin{proof}
    \ref{itm:prop:Icomparandenvelope:1} $\implies$ \ref{itm:prop:Icomparandenvelope:2}. This follows immediately from \cref{cor:Ienvelopedef2}.

    \ref{itm:prop:Icomparandenvelope:2} $\implies$ \ref{itm:prop:Icomparandenvelope:1}. It follows from \cref{prop:Ienvprojection} that
    \[
        \varphi\sim_{\mathcal{I}} P_{\theta}[\varphi]_{\mathcal{I}}\leq P_{\theta}[\psi]_{\mathcal{I}} \sim_{\mathcal{I}} \psi.
    \]
    We conclude \ref{itm:prop:Icomparandenvelope:1}.
\end{proof}

\begin{example}\label{ex:Pequivnotequivex}
    Let us continue our example \cref{ex:loglogsing}, where $X=\mathbb{P}^1$, $\omega$ is the Fubini--Study metric and $\varphi\in \PSH(X,\omega)$ has log-log singularity at $0$. We have shown that $P_{\omega}[\varphi]=0$ in \eqref{eq:Pomegavarphi0temp1}, so $\varphi\sim_{P} 0$ and hence $\varphi\sim_{\mathcal{I}} 0$. In particular, $P$-equivalence is not equivalent to the equivalence of singularity types.

    On the other hand, consider a potential $\psi\in \PSH(X,\omega)$ with log singularity at $0$, as in \cref{ex:logsingP1}. We know that $\nu(\psi,0)=1$ from the explicit expression \eqref{eq:logsingP1ex}. So $\psi\not\sim_{\mathcal{I}} 0$ and hence $\psi\not\sim_{P} 0$.

    Moreover, $\psi \preceq_{P} \varphi$ and hence $\psi\preceq_{\mathcal{I}}\varphi$.
\end{example}


We give an example showing that $P$-equivalence is not equivalent to $\mathcal{I}$-equivalence.
\begin{example}\label{ex:BBJ}
    Let $X=\mathbb{P}^1$ and $\omega$ be the Fubini--Study metric. Let $K\subseteq \mathbb{P}^1$ be a polar Cantor set carrying an atom free probability measure $\mu$ supported on $K$ (see \cite[Page~31]{Car83}). Write $\mu=\omega+\ddc\varphi$ for some $\omega$-subharmonic function $\varphi$. Since $\mu$ is atom free, we know that all Lelong numbers of $\varphi$ are $0$. On the other hand, $\varphi$ has $0$ non-pluripolar mass since $K$ is pluripolar.

    Then observe that $\varphi\sim_{\mathcal{I}} 0$ while $\varphi\not\sim_P 0$.
\end{example}

For later use, we give the following definition.
\begin{definition}\label{def:elemetric}
    Let $L$ be a pseudo-effective line bundle on $X$. An \emph{elementary metric}\index{elementary metric} on $L$ is a psh metric $h$ on $L$ such that there is a generalized Fubini--Study metric $h'$ on $L$ such that
    \[
    \ddc h\sim_{P} \ddc h'.
    \]
    The set of elementary metrics on $L$ is denoted by $\Ele(L)$\index{$\Ele(L)$}.

    We also say $\ddc h$ is elementary. If we have fixed a Hermitian metric $h_0$ on $L$, and if we represent $h$ as $h_0\exp(-\varphi)$, we also say the quasi-psh function $\varphi$ is elementary.
\end{definition}
Recall that the generalized Fubini--Study metrics are defined in \cref{def:gFSmetric}.

\subsection{Properties of the preorders}
\label{subsec:preorder_prop}

Now we state a more natural version of the monotonicity theorem \cref{thm:mono}.
\begin{proposition}[Monotonicity theorem]\label{prop:mono2}
    Let $\theta_1,\ldots,\theta_n$ be closed real smooth $(1,1)$-forms on $X$.
    Let $\varphi_i,\psi_i\in \PSH(X,\theta_i)$ for $i=1,\ldots,n$. Assume that $\varphi_i\preceq_P \psi_i$ for each $i$. Then
    \[
        \int_X \theta_{1,\varphi_1}\wedge \cdots \wedge  \theta_{n,\varphi_n}\leq \int_X \theta_{1,\psi_1}\wedge \cdots \wedge  \theta_{n,\psi_n}.
    \]
    In other words, the non-pluripolar mass is a decreasing quantity with respect to the $P$-singularity types of the potentials.
\end{proposition}
\begin{proof}
    Fix a Kähler form $\omega$ on $X$. For each $i=1,\ldots,n$, since $\varphi_i\preceq_P\psi_i$, we have
    \[
        P_{\theta_i+\epsilon\omega}[\varphi_i]\leq P_{\theta_i+\epsilon\omega}[\psi_i]
    \]
    for all $\epsilon>0$. Therefore, by \cref{prop:Ppresmass} and our previous monotonicity theorem \cref{thm:mono}, we have
    \[
        \int_X (\theta_1+\epsilon\omega)_{\varphi_1}\wedge \cdots \wedge (\theta_n+\epsilon\omega)_{\varphi_n}\leq \int_X (\theta_1+\epsilon\omega)_{\psi_1}\wedge \cdots \wedge (\theta_n+\epsilon\omega)_{\psi_n}.
    \]
    Letting $\epsilon\to 0+$ and applying the multi-linearity of the non-pluripolar products \itemref{itm:prop:npp1:6}, we find the desired inequality.
\end{proof}

The $P$-preorder allows us to derive a natural generalization of \cref{lma:pathoenvelope}.
\begin{theorem}\label{thm:pathoenvelope}
    Let $\varphi,\psi\in \PSH(X,\theta)_{>0}$, $\varphi\preceq_P \psi$. Then for any
    \begin{equation}\label{eq:arangetemp_reproduce}
        a\in \left( 1, \left(\frac{\int_X \theta_{\psi}^n}{\int_X \theta_{\psi}^n-\int_X \theta_{\varphi}^n}\right)^{1/n} \right),
    \end{equation}
    the function
    \[
        \eta\coloneqq P_{\theta}\Bigl(a\varphi+(1-a)\psi\Bigr)
    \]
    lies in $\PSH(X,\theta)_{>0}$, satisfies $\eta\preceq_P \varphi$ and
    \begin{equation}\label{eq:path_env_reproduce}
        a^{-1}\eta+\left(1-a^{-1}\right)\psi\leq \varphi.
    \end{equation}

    If in addition, $\varphi\sim_P \psi$, then $\eta\sim_P \varphi$ as well.
\end{theorem}
In \eqref{eq:arangetemp_reproduce}, the fraction is understood as $\infty$ if either $\int_X \theta_{\psi}^n=\int_X \theta_{\varphi}^n$ or $n=0$, exactly as in \cref{lma:pathoenvelope}.
\begin{proof}
    By our assumption and \eqref{eq:varphipreceqP}, we have $\varphi\preceq P_{\theta}[\varphi]\leq P_{\theta}[\psi]$.
    Due to \cref{prop:Ppresmass}, the interval \eqref{eq:arangetemp_reproduce} remains unchanged after replacing $\psi$ by $P_{\theta}[\psi]$. Fix $a$ in this interval.
    So by \cref{lma:pathoenvelope} and the discussions afterwards,
    \[
        \eta'\coloneqq  P_{\theta}\Bigl(a\varphi+(1-a)P_{\theta}[\psi]\Bigr)\in \PSH(X,\theta)_{>0}.
    \]
    We first check that $\eta$ is a genuine $\theta$-psh function. Choose a closed non-pluripolar set $K\subseteq X$ on which $\varphi$ and $\psi$ are bounded. Then the obstacle $a\varphi+(1-a)\psi$ is bounded above on $K$.
    Every candidate in the definition of $\eta$ is therefore uniformly bounded above on $X$ by \cref{rmk:L1cpt}, so $\eta\not\equiv \infty$ and hence $\eta\in \PSH(X,\theta)\cup\{-\infty\}$.
    On the other hand, since $\psi\preceq P_{\theta}[\psi]$, we can find $C>0$ such that $\psi\leq P_{\theta}[\psi]+C$. Then
    \[
        \eta'-(a-1)C\leq a\varphi+(1-a)\psi
    \]
    quasi-everywhere, hence $\eta\geq \eta'-(a-1)C$. Thus $\eta'\preceq \eta$, and the monotonicity theorem \cref{thm:mono} gives $\eta\in \PSH(X,\theta)_{>0}$.

    By \eqref{eq:Pthetafleqfqe},
    \[
        \eta\leq a\varphi+(1-a)\psi
    \]
    quasi-everywhere. Equivalently, \eqref{eq:path_env_reproduce} holds quasi-everywhere. Since both sides of \eqref{eq:path_env_reproduce} are $\theta$-psh, \cref{prop:pshfuncdetdense} upgrades the inequality to an everywhere inequality.

    It remains to prove the assertion about the $P$-order.
    Using \cref{prop:Pconc}, we get
    \[
        a^{-1} P_{\theta}[\eta]+\left(1-a^{-1}\right)P_{\theta}[\psi] \leq P_{\theta}[\varphi].
    \]
    By the assumption $\varphi\preceq_P\psi$, we also have $P_{\theta}[\varphi]\leq P_{\theta}[\psi]$. Therefore,
    \[
        a^{-1} P_{\theta}[\eta]\leq P_{\theta}[\varphi]-\left(1-a^{-1}\right)P_{\theta}[\psi]\leq a^{-1} P_{\theta}[\varphi]
    \]
    quasi-everywhere and hence everywhere by \cref{prop:pshfuncdetdense}. It follows that $P_{\theta}[\eta]\leq P_{\theta}[\varphi]$, namely $\eta\preceq_P\varphi$.

    Finally assume that $\varphi\sim_P\psi$. Then $\int_X\theta_{P_{\theta}[\psi]}^n=\int_X\theta_{\varphi}^n$ by \cref{prop:Ppresmass}.
    Applying \cref{cor:pathoenvelopeeqmass} gives
    \[
        \int_X\theta_{\eta'}^n=\int_X\theta_{\varphi}^n.
    \]
    Since $\eta'\preceq\eta$, the monotonicity theorem \cref{thm:mono} gives
    \[
        \int_X\theta_{\eta}^n\geq \int_X\theta_{\varphi}^n.
    \]
    The reverse inequality follows from $\eta\preceq_P\varphi$ and \cref{prop:mono2}. Hence $\eta\sim_P\varphi$ by \cref{prop:Pequivchar2}.
\end{proof}


Next we show that the $P$ and $\mathcal{I}$-preorders are preserved by some natural operations.

\begin{lemma}\label{lma:Ipartorderinv}
Let $\pi\colon Y\rightarrow X$ be a proper bimeromorphic morphism from a Kähler manifold $Y$.
    Given two quasi-plurisubharmonic functions $\varphi,\psi$ on $X$, then the following are equivalent:
    \begin{enumerate}
        \item\label{itm:lma:Ipartorderinv:1} $\varphi\preceq_{P} \psi$;
        \item\label{itm:lma:Ipartorderinv:2} $\pi^*\varphi\preceq_{P} \pi^*\psi$.
    \end{enumerate}
    The same holds with $\mathcal{I}$ in place of $P$.
\end{lemma}
\begin{proof}
   In the $P$-case, this follows from \cref{prop:Penvbimero}, while in the $\mathcal{I}$-case, this follows from \cref{prop:Ienvelopebimero}.
\end{proof}

\begin{proposition}\label{prop:Ppartialsum}
    Let $\varphi,\psi,\varphi',\psi'\in \QPSH(X)$. Assume that
    \[
    \varphi\preceq_P \psi,\quad \varphi'\preceq_P \psi'.
    \]
    Then
    \[
    \varphi+\varphi'\preceq_P \psi+\psi'.
    \]
    The same holds with $\preceq_{\mathcal{I}}$ in place of $\preceq_P$.
\end{proposition}
\begin{proof}
    Take a Kähler form $\omega$ on $X$ such that $\varphi,\psi,\varphi',\psi'\in \PSH(X,\omega)_{>0}$. The statement for $\preceq_{\mathcal{I}}$ is a simple consequence of \cref{prop:Lelongmax}. We only need to handle the case of $\preceq_{P}$.

    \textbf{Step~1}.
    We first show that
    \[
    P_{\omega}[\varphi]+P_{\omega}[\varphi']\sim_P \varphi+\varphi'.
    \]
    In fact, we clearly have
    \[
    P_{\omega}[\varphi]+P_{\omega}[\varphi']\succeq \varphi+\varphi'.
    \]
    So by \cref{prop:Pequivchar2}, it suffices to show that they have the same mass. We compute
    \[
    \begin{aligned}
    &\int_X \Bigl(2\omega+\ddc P_{\omega}[\varphi]+\ddc P_{\omega}[\varphi']\Bigr)^n\\
    =& \sum_{j=0}^n \binom{n}{j} \int_X \Bigl(\omega+\ddc P_{\omega}[\varphi]\Bigr)^j\wedge \Bigl(\omega+\ddc P_{\omega}[\varphi']\Bigr)^{n-j}\\
    =& \sum_{j=0}^n \binom{n}{j} \int_X \omega_{\varphi}^j\wedge \omega_{\varphi'}^{n-j}\\
    =& \int_X \left(2\omega+\ddc\varphi+\ddc\varphi'\right)^n,
    \end{aligned}
    \]
    where we applied \cref{prop:Ppresmass} to the mixed products on the third line.

    \textbf{Step~2}. By Step~1, we may assume that $\varphi,\psi,\varphi',\psi'$ are all model potentials, after replacing them by their $P_{\omega}$-envelopes. So $\varphi\leq \psi$ and $\varphi'\leq \psi'$. Our assertion follows.
\end{proof}
\begin{proposition}\label{prop:Ppartialsup}
    Let $(\varphi_i)_{i\in I}$, $(\psi_i)_{i\in I}$ be uniformly bounded from above non-empty families in $\QPSH(X)$. Assume that there exists a closed smooth real $(1,1)$-form $\theta$ such that $\varphi_i,\psi_i\in \PSH(X,\theta)$ and $\varphi_i\preceq_P \psi_i$ for all $i\in I$. Then
    \[
    \uscsup[i\in I] \varphi_i\preceq_P \uscsup[i\in I]\psi_i.
    \]

    The same holds with $\preceq_{\mathcal{I}}$ in place of $\preceq_P$.
\end{proposition}
\begin{proof}
    By adding a Kähler form to $\theta$, we may assume that $\varphi_i,\psi_i\in \PSH(X,\theta)_{>0}$ for all $i\in I$. The statement for $\preceq_{\mathcal{I}}$ is a simple consequence of \cref{cor:supsLelong}, we only have to consider the statement for $\preceq_P$.

    \textbf{Step~1}. We first handle the case where $I$ is a directed set and $(\varphi_i)_{i\in I}$ and $(\psi_i)_{i\in I}$ are increasing nets.

   In this case,  our assertion follows simply from \cref{prop:incnetmodel}.

    \textbf{Step~2}. We handle the case where $I$ is finite. We may assume that $I=\{0,1\}$.
    It suffices to show that
    \[
        P_{\theta}[\varphi_0]\lor P_{\theta}[\varphi_1]\sim_P \varphi_0\lor\varphi_1,
    \]
    which follows from \cref{lma:Plor1}.

    \textbf{Step~3}. The general case can be reduced to the two cases handled in Step~1 and Step~2. More precisely,
    by applying \cref{prop:Choquet} to the two families and taking the union of the resulting countable subsets, we can find a countable subset $J\subseteq I$ such that
    \[
    \uscsup[j\in J] \varphi_j=\uscsup[i\in I] \varphi_i,\quad \uscsup[j\in J] \psi_j=\uscsup[i\in I] \psi_i.
    \]
    We may replace $I$ by $J$ and assume that $I$ is countable. If $I$ is finite, Step~2 applies directly. Otherwise, write $I=\mathbb{Z}_{>0}$ and set
    \[
        \tilde{\varphi}_m=\varphi_1\lor\cdots\lor\varphi_m,\quad
        \tilde{\psi}_m=\psi_1\lor\cdots\lor\psi_m.
    \]
    By Step~2, $\tilde{\varphi}_m\preceq_P\tilde{\psi}_m$ for every $m$. The two new sequences are increasing and have the same upper envelopes as the original families, so Step~1 applies.

\end{proof}

We now refine \cref{lma:landcorrectmass} using the $P$-equivalence relation.
\begin{proposition}\label{prop:rooftopprePequiv}
    Let $\varphi,\psi,\varphi',\psi'\in \PSH(X,\theta)$ for some closed smooth real $(1,1)$-form $\theta$ on $X$. Assume that
    \[
    \varphi\sim_P \varphi',\quad \psi\sim_P\psi',\quad  \varphi'\land \psi'\in \PSH(X,\theta)_{>0}.
    \]
    Then
    \[
    \varphi\land \psi\in \PSH(X,\theta)_{>0},\quad \varphi\land \psi\sim_P \varphi'\land \psi'.
    \]
\end{proposition}
\begin{proof}
    We first observe that $\varphi,\varphi',\psi,\psi'\in \PSH(X,\theta)_{>0}$ by assumption. Let
    \[
    \phi\coloneqq P_{\theta}[\varphi]=P_{\theta}[\varphi'],\quad \gamma=P_{\theta}[\psi]=P_{\theta}[\psi'].
    \]
    Since $\varphi'\preceq\phi$ and $\psi'\preceq\gamma$, after adding a common constant to $\varphi'$ and $\psi'$, we may assume that $\varphi'\leq\phi$ and $\psi'\leq\gamma$. Then $\varphi'\land\psi'\leq\phi\land\gamma$, so the assumption $\varphi'\land\psi'\in \PSH(X,\theta)_{>0}$ implies $\phi\land \gamma \in \PSH(X,\theta)_{>0}$ by monotonicity. It follows from \cref{lma:landcorrectmass} that
    \[
    \int_X \theta_{\varphi'\land \psi'}^n=\int_X \theta_{\phi\land \gamma}^n.
    \]
    Next, we apply \cref{lma:landcorrectmass} again to conclude that $\varphi\land \psi\in \PSH(X,\theta)$ and
    \[
    \int_X \theta_{\varphi\land \psi}^n=\int_X \theta_{\phi\land \gamma}^n=\int_X \theta_{\varphi'\land \psi'}^n>0.
    \]
    Since $\varphi\land\psi\preceq \phi\land\gamma$ and the two potentials have the same positive mass, \cref{thm:Pvarphidiffdef} gives $P_{\theta}[\varphi\land\psi]=\phi\land\gamma$. The same argument gives $P_{\theta}[\varphi'\land\psi']=\phi\land\gamma$, hence $\varphi\land\psi\sim_P\varphi'\land\psi'$.
\end{proof}

We can formulate a generalization of \cref{prop:envrelfullmass} using the $P$-partial order.
\begin{corollary}\label{cor:P_par_envelope_rel}
    Let $\varphi,\psi\in \PSH(X,\theta)_{>0}$. Assume that $\psi\preceq_P \varphi$. Then for any $C>0$, $\psi\land (\varphi+C)\in \PSH(X,\theta)_{>0}$ and
    \begin{equation}\label{eq:supvarphipsipC}
        P_{\theta}[\varphi](\psi)=\uscsup[C>0] \psi\land (\varphi+C)=\psi.
    \end{equation}
\end{corollary}
We refer to \cref{def:Pthetavarphif} for the definition of $P_{\theta}[\varphi](\psi)$.
\begin{proof}
    Since $\psi\preceq_P\varphi$, by \cref{prop:rooftopprePequiv}, we get
    \[
        \psi\land(\varphi+C)\in \PSH(X,\theta)_{>0}
    \]
    and $\psi\land(\varphi+C)\sim_P\psi$ for every $C>0$.

    Let
    \[
        \eta\coloneqq \uscsup[C>0]\psi\land(\varphi+C).
    \]
    Then $\eta\sim_P \psi$.
    By \cref{cor:Pvarphisupport}, we have
    \[
        \theta_\eta^n\leq \mathds{1}_{\{\eta=\psi\}}\,\theta_\psi^n.
    \]
    The domination principle \cref{thm:dom_prin_1} gives $\eta\geq \psi$. The reverse inequality is trivial, so $\eta=\psi$.
\end{proof}


The $P$-equivalence relation characterizes when a subgeodesic exists. See \cref{def:subgeod} for the notion of subgeodesics.
\begin{theorem}\label{thm:geodexi}
    Let $\varphi_0,\varphi_1\in \PSH(X,\theta)_{>0}$. Then the following are equivalent:
    \begin{enumerate}
        \item\label{itm:thm:geodexi:1} There is a subgeodesic from $\varphi_0$ to $\varphi_1$;
        \item\label{itm:thm:geodexi:2} $\varphi_0\sim_P\varphi_1$.
    \end{enumerate}
\end{theorem}
\begin{proof}
    \ref{itm:thm:geodexi:2} $\implies$ \ref{itm:thm:geodexi:1}. This follows from \cref{prop:perronenvissubgeod}.

    \ref{itm:thm:geodexi:1} $\implies$ \ref{itm:thm:geodexi:2}.
    Let $(\varphi_t)_{t\in (0,1)}$ be a subgeodesic from $\varphi_0$ to $\varphi_1$.

    \textbf{Step~1}. We assume that $\varphi_0\geq \varphi_1$.

    We first replace the given subgeodesic by the geodesic from $\varphi_0$ to $\varphi_1$. Note that the geodesic is necessarily a subgeodesic as well by \cref{lma:geodsubgeod}.

    Observe that $t\mapsto\varphi_t$ is decreasing due to \cref{lma:geoddecreasing}.

    Let $\varphi_t=\varphi_1$ for all $t>1$. Then by the gluing lemma \cref{lma:subgeodgluing}, we find that $(\varphi_t)_{t\geq 0}$ is a subgeodesic ray.

    Next, we consider the Legendre transform
    \[
    \Gamma_{\tau}\coloneqq \inf_{t\geq 0}(\varphi_t-t\tau),\quad \tau\in \mathbb{R}.
    \]
    It follows from Kiselman's principle \cref{prop:Kis} that $\Gamma_{\tau}\in \PSH(X,\theta)\cup \{-\infty\}$. Note that for $\tau>0$, we clearly have $\Gamma_{\tau}\equiv -\infty$. On the other hand, for $\tau\leq 0$,
    \[
    \Gamma_{\tau}=\inf_{t\in [0,1]}(\varphi_t-t\tau)\in \PSH(X,\theta).
    \]

    By Legendre inversion (cf. \cref{cor:Leginvpartdef}), for $t>0$,
    \[
    \varphi_t=\sup_{\tau<0} \left(\Gamma_{\tau}+t\tau\right).
    \]
    Fix a K\"ahler form $\omega$ on $X$.
    It follows from \cref{prop:Ppartialsup} that for each $t>0$,
    \[
    \varphi_t\sim_P \uscsup[\tau<0] P_{\theta+\omega}[\Gamma_{\tau}].
    \]
    The right-hand side is independent of $t$.
    Here by adding $\omega$, we no longer have to worry about the possibility where $\Gamma_{\tau}$ has vanishing mass.

    Write
    \[
    \varphi_0=\uscsup[t\in (0,1)]\varphi_t.
    \]
    By \cref{prop:Ppartialsup} again, we find that
    \[
    \varphi_0\sim_P \uscsup[\tau<0] P_{\theta+\omega}[\Gamma_{\tau}]
    \]
    as well. So $\varphi_0\sim_P \varphi_1$.


    \textbf{Step~2}. We prove the general case.

    Observe that $(\varphi_t\lor\varphi_1)_{t\in (0,1)}$ is a subgeodesic from $\varphi_0\lor \varphi_1$ to $\varphi_1$. By Step~1,
    \[
    \varphi_0\lor \varphi_1\sim_P \varphi_1.
    \]
    Hence $\varphi_0\preceq_P \varphi_1$ by \cref{lma:reform_preceqP}. The converse is proved similarly. Hence \ref{itm:thm:geodexi:2} follows.

\end{proof}

Restating the \ref{itm:thm:geodexi:1} $\implies$ \ref{itm:thm:geodexi:2} direction in the theorem in terms of the complexification, we find the following interesting result.

Let $S=\{z\in \mathbb{C}:0< \Real z <1\}$. We write $p_1\colon X\times S\rightarrow X$ for the natural projection.
\begin{corollary}\label{cor:comp:L653}
    Let $\Phi\in \PSH(X\times S,p_1^*\theta)$. Assume that for any $c\in \mathbb{R}$, $x\in X$ and $z\in S$, we have
    \[
    \Phi(x,z)=\Phi(x,z+\mathrm{i}c).
    \]
    Then $\int_X (\theta+\ddc \Phi_z)^n$ is independent of $z\in S$, where $\Phi_z \in \PSH(X,\theta)$ is given by $\Phi_z(x)=\Phi(x,z)$.
\end{corollary}

\begin{proof}
    Fix a Kähler form $\omega$ on $X$. We may regard $\Phi$ as an element of $\PSH(X\times S,p_1^*(\theta+\omega))$. For any $0<a<b<1$, the restriction $(\Phi_t)_{t\in [a,b]}$ is then a subgeodesic from $\Phi_a$ to $\Phi_b$ in the class $\theta+\omega$. The slices have positive mass in this class, so \cref{thm:geodexi} gives $\Phi_a\sim_P\Phi_b$. By \cref{prop:mono2}, this directly implies
    \[
        \int_X \theta_{\Phi_a}^n=\int_X \theta_{\Phi_b}^n.
    \]
    Since $\Phi$ is invariant under imaginary translations, $\Phi_z$ depends only on $\Real z$, so the mass is independent of $z\in S$.
\end{proof}

\section[The \texorpdfstring{$\lowercase{d}_S$}{dS}-pseudometric]{The \texorpdfstring{$d_S$}{dS}-pseudometric}\label{sec:dsdef}
Let $X$ be a connected compact Kähler manifold of dimension $n$ and $\theta$ be a closed real smooth $(1,1)$-form on $X$ representing a big cohomology class. The goal of this section is to study a pseudometric on the space $\PSH(X,\theta)$.


\subsection{The definition of the \texorpdfstring{$d_S$}{dS}-pseudometric}
\label{subsec:dS_def}
Recall that for any $\varphi\in \PSH(X,\theta)$, the geodesic ray $\ell^{\varphi}\in \mathcal{R}^1(X,\theta)$ is defined in \cref{ex:rayasspsh}.
\begin{definition}\label{def:dS}
    For $\varphi,\psi\in \PSH(X,\theta)$, we define
    \[
        d_S(\varphi,\psi)\coloneqq d_1(\ell^{\varphi},\ell^{\psi}).
    \]
    \index{$d_S$}
    When we want to be more specific, we write $d_{S,\theta}$ instead of $d_S$.
\end{definition}
The $d_1$ distance of geodesic rays is defined in \cref{def:d1rays}.

\begin{proposition}\label{prop:comp:L686}
    The function $d_S$ defined in \cref{def:dS} is a pseudometric on $\PSH(X,\theta)$.
\end{proposition}
\begin{proof}
    This follows immediately from \cref{thm:d1raycomplete}.
\end{proof}
When studying a pseudometric, the first thing is to understand when the distance between two elements vanishes.

We first prove a preparation:
\begin{lemma}\label{lma:dSalmostriang}
    Let $\varphi,\psi\in \PSH(X,\theta)$. Then
    \[
        d_S(\varphi,\psi)\leq d_S(\varphi,\varphi\lor \psi)+d_S(\psi,\varphi\lor \psi) \leq C_n d_S(\varphi,\psi),
    \]
    where $C_n=3(n+1)2^{n+2}$.
\end{lemma}
We shall use the notations introduced in \cref{ex:rayasspsh}.
\begin{proof}
    We claim that
    \begin{equation}\label{eq:elllorsingtype}
        \ell^{\varphi}\lor \ell^{\psi}=\ell^{\varphi\lor \psi}.
    \end{equation}
    Recall that $\lor$ is defined in \cref{def:lorray1}.
    Note that this assertion implies our desired inequality by \cref{lma:d1rayineq}.

    In proving this assertion, we may assume that $\varphi,\psi\leq 0$ since
    \[
    \ell^{\varphi+C}=\ell^{\varphi},\quad \ell^{\psi+C}=\ell^{\psi},\quad \ell^{(\varphi+C)\lor (\psi+C)}=\ell^{\varphi\lor \psi}
    \]
    for any $C\in \mathbb{R}$.

    In fact, it is clear that
    \[
        \ell^{\varphi}\leq   \ell^{\varphi\lor \psi},\quad \ell^{\psi}\leq   \ell^{\varphi\lor \psi},
    \]
    so the $\leq$ direction in \eqref{eq:elllorsingtype} holds.

    Conversely, if $\ell'\in \mathcal{R}^1(X,\theta)$ and $\ell'\geq \ell^{\varphi}\lor \ell^{\psi}$, then for each $t\geq 0$,
    \[
        \ell'_t\geq \Bigl((V_{\theta}-t)\lor \varphi\Bigr) \lor \Bigl((V_{\theta}-t)\lor \psi\Bigr)=(V_{\theta}-t)\lor (\varphi\lor \psi).
    \]
    See \eqref{eq:ellphitexp} for the first inequality.
    Therefore, thanks to the definition of the Perron envelopes (cf. \eqref{eq:Perron2}), we have
    \[
    \ell'_s\geq \ell^{\varphi\lor \psi,t}_s
    \]
    for any $0\leq s\leq t$.
    It follows from \eqref{eq:ellphitexp} that $\ell'_s\geq \ell^{\varphi\lor \psi}_s$ for any $s\geq 0$.
\end{proof}


\begin{proposition}\label{prop:ds0char}
    Let $\varphi,\psi\in \PSH(X,\theta)$. Then the following are equivalent:
    \begin{enumerate}
        \item\label{itm:prop:ds0char:1} $\varphi\sim_P \psi$;
        \item\label{itm:prop:ds0char:2} $d_S(\varphi,\psi)=0$.
    \end{enumerate}

    In particular, $d_S(\varphi,P_{\theta}[\varphi])=0$ for all $\varphi\in \PSH(X,\theta)_{>0}$.
\end{proposition}
\begin{proof}
    By \cref{lma:reform_preceqP}, we have $\varphi\sim_P \psi$ if and only if $\varphi\sim_P \varphi\lor \psi$ and $\psi\sim_P \varphi\lor \psi$. By \cref{lma:dSalmostriang}, $d_S(\varphi,\psi)=0$ if and only if $d_S(\varphi,\varphi\lor \psi)=0$ and $d_S(\psi,\varphi\lor \psi)=0$. So it suffices to prove the assertion when $\varphi\leq \psi$. Assuming this, by \cref{prop:d1geod_diff_E}, \ref{itm:prop:ds0char:2} holds if and only if
    \[
        \mathbf{E}\left(\ell^{\varphi}\right)=  \mathbf{E}\left(\ell^{\psi}\right),
    \]
    where $\mathbf{E}$ is introduced in \cref{def:radialMAenergy2}.
    But by \eqref{eq:Elphi}, this holds if and only if
    \[
        \sum_{j=0}^n \int_X \theta_{\varphi}^j\wedge \theta_{V_{\theta}}^{n-j}=\sum_{j=0}^n \int_X \theta_{\psi}^j\wedge \theta_{V_{\theta}}^{n-j}.
    \]
    Thanks to the monotonicity theorem \cref{thm:mono}, this holds if and only if for all $j=0,\ldots,n$,
    \[
        \int_X \theta_{\varphi}^j\wedge \theta_{V_{\theta}}^{n-j}=\int_X \theta_{\psi}^j\wedge \theta_{V_{\theta}}^{n-j},
    \]
    which is equivalent to \ref{itm:prop:ds0char:1} by \cref{prop:Pequivchar2}.
\end{proof}


\begin{lemma}\label{lma:varphileqpsi_metric}
    Suppose that $\varphi,\psi\in \PSH(X,\theta)$ and $\varphi\preceq_P \psi$. Then
    \[
    d_S(\varphi,\psi)=\frac{1}{n+1}\sum_{j=0}^n \left(\int_X \theta_{\psi}^j\wedge \theta_{V_{\theta}}^{n-j}-\int_X \theta_{\varphi}^j\wedge \theta_{V_{\theta}}^{n-j}\right).
    \]
\end{lemma}
\begin{proof}
    Set $\chi\coloneqq\varphi\lor\psi$. By \cref{lma:reform_preceqP}, we have $\chi\sim_P\psi$. Hence \cref{prop:ds0char} gives $d_S(\chi,\psi)=0$, and the pseudometric property gives
    \[
        d_S(\varphi,\psi)=d_S(\varphi,\chi).
    \]
    Since $\varphi\leq\chi$, \eqref{eq:Elphi} and \cref{prop:d1geod_diff_E} give
    \[
    d_S(\varphi,\chi)=\frac{1}{n+1}\sum_{j=0}^n \left(\int_X \theta_{\chi}^j\wedge \theta_{V_{\theta}}^{n-j}-\int_X \theta_{\varphi}^j\wedge \theta_{V_{\theta}}^{n-j}\right).
    \]
    Finally, \cref{prop:Pequivchar2} applied to $\psi\leq\chi$ gives
    \[
        \int_X \theta_{\chi}^j\wedge \theta_{V_{\theta}}^{n-j}=\int_X \theta_{\psi}^j\wedge \theta_{V_{\theta}}^{n-j}
    \]
    for all $j=0,\ldots,n$.
\end{proof}
\begin{corollary}\label{cor:dsthreeterm}
    Suppose that $\varphi,\psi,\eta\in \PSH(X,\theta)$ and $\varphi\preceq_P\psi\preceq_P \eta$. Then
    \[
        d_S(\varphi,\eta)\geq d_S(\varphi,\psi),\quad d_S(\varphi,\eta)\geq d_S(\psi,\eta).
    \]
\end{corollary}
\begin{proof}
    This is an immediate consequence of \cref{lma:varphileqpsi_metric} and the monotonicity theorem \cref{prop:mono2}.
\end{proof}

\begin{corollary}\label{cor:dsmetricdoubleineq}
    For any $\varphi,\psi\in \PSH(X,\theta)$, we have
    \begin{equation}\label{eq:ds_biineq}
        \begin{split}
            d_S(\varphi,\psi)\leq & \frac{1}{n+1}\sum_{j=0}^n \left( 2\int_X \theta_{\varphi\lor \psi}^j\wedge \theta_{V_{\theta}}^{n-j}-\int_X \theta_{\varphi}^j\wedge \theta_{V_{\theta}}^{n-j}-\int_X \theta_{\psi}^j\wedge \theta_{V_{\theta}}^{n-j} \right) \\
            \leq & C_n d_S(\varphi,\psi),
        \end{split}
\end{equation}
where $C_n=3(n+1)2^{n+2}$.

In particular, if $(\varphi_i)_{i\in I}$ is a net in $\PSH(X,\theta)$ with $d_S$-limit $\varphi$, then for each $j=0,\ldots,n$,
\[
    \lim_{i\in I}\int_X \theta_{\varphi_i}^j\wedge \theta_{V_{\theta}}^{n-j}=\int_X \theta_{\varphi}^j\wedge \theta_{V_{\theta}}^{n-j}=\lim_{i\in I}\int_X \theta_{\varphi_i\lor \varphi}^j\wedge \theta_{V_{\theta}}^{n-j}.
\]
\end{corollary}
\begin{proof}
    The estimates \eqref{eq:ds_biineq} follow from the combination of \cref{lma:varphileqpsi_metric} and \cref{lma:dSalmostriang}.

    Suppose that $\varphi_i\xrightarrow{d_S}\varphi$. Then $\varphi_i\lor \varphi\xrightarrow{d_S}\varphi$ by \cref{lma:dSalmostriang}. Therefore, \cref{thm:mono} and \cref{lma:varphileqpsi_metric} imply that
    \[
    \lim_{i\in I}\int_X \theta_{\varphi_i\lor \varphi}^j\wedge \theta_{V_{\theta}}^{n-j}= \int_X \theta_{\varphi}^j\wedge \theta_{V_{\theta}}^{n-j}
    \]
    for any $j=0,\ldots,n$.
    The last assertion now follows from \eqref{eq:ds_biineq} and \cref{thm:mono}.
\end{proof}

\begin{corollary}\label{cor:dSconv_bimero}
    Let $\pi\colon Y\rightarrow X$ be a proper bimeromorphic morphism from a Kähler manifold. Given a net $(\varphi_i)_{i\in I}$ in $\PSH(X,\theta)$ and $\varphi\in \PSH(X,\theta)$, the following are equivalent:
    \begin{enumerate}
        \item $\varphi_i\xrightarrow{d_S}\varphi$;
        \item $\pi^*\varphi_i\xrightarrow{d_S}\pi^*\varphi$.
    \end{enumerate}
\end{corollary}
\begin{proof}
    It suffices to observe that the middle term in \eqref{eq:ds_biineq} is invariant under $\pi^*$.
\end{proof}

\begin{corollary}\label{cor:incnetdSconv}
    Suppose that $(\varphi_i)_{i\in I}$ is an increasing net in $\PSH(X,\theta)$, uniformly bounded from above. Then
    \[
        \varphi_i\xrightarrow{d_S} \uscsup[i\in I]\varphi_i.
    \]
\end{corollary}
If the $\varphi_i$'s are all model potentials in $\PSH(X,\theta)_{>0}$, then so is $\uscsup[i\in I]\varphi_i$, as we have seen in \cref{prop:incnetmodel}.
\begin{proof}
    Write $\varphi=\uscsup[j\in I]\varphi_j$. Recall that by \cref{prop:pshfunction_closedseq}, $\varphi\in \PSH(X,\theta)$.
    By \cref{lma:varphileqpsi_metric}, it suffices to show that for each $k=0,\ldots,n$, we have
    \[
        \lim_{j\in I}\int_X \theta_{\varphi_j}^k\wedge \theta_{V_{\theta}}^{n-k} =  \int_X \theta_{\varphi}^k\wedge \theta_{V_{\theta}}^{n-k}.
    \]
    The latter follows from \cref{cor:incseqnppcont}.
\end{proof}

\begin{corollary}\label{cor:massdiffdS}
    Let $\varphi,\psi\in \PSH(X,\theta)$. Then
    \[
    \left| \int_X \theta_{\varphi}^n-\int_X \theta_{\psi}^n \right|\leq D_n d_S(\varphi,\psi),
    \]
    where $D_n=3(n+1)C_n$ with $C_n$ being the same constant as in \cref{lma:dSalmostriang}.
\end{corollary}
In other words, the non-pluripolar mass is uniformly continuous with respect to the $d_S$-pseudometric.
\begin{proof}
    We compute
    \[
    \begin{aligned}
        \left| \int_X \theta_{\varphi}^n-\int_X \theta_{\psi}^n \right|\leq & \left|2\int_X \theta_{\varphi\lor \psi}^n-\int_X \theta_{\varphi}^n-\int_X \theta_{\psi}^n\right| + 2\left| \int_X \theta_{\varphi\lor \psi}^n-\int_X \theta_{\varphi}^n\right|\\
        \leq &  (n+1)C_n d_S(\varphi,\psi)+2(n+1) d_S(\varphi,\varphi\lor \psi)\\
        \leq & (n+1)C_n d_S(\varphi,\psi)+2(n+1)C_n d_S(\varphi,\psi),
    \end{aligned}
    \]
    where the first line is just the triangle inequality, the second line follows from \cref{cor:dsmetricdoubleineq} and the third line follows from \cref{lma:dSalmostriang}.
\end{proof}

By contrast to the increasing case, for nets decreasing with respect to the $P$-preorder, the situation is different:
\begin{corollary}\label{cor:decnetdS}
    Suppose that $(\varphi_i)_{i\in I}$ is a net in $\PSH(X,\theta)$ such that
    \[
        \varphi_j\preceq_P\varphi_i
    \]
    whenever $j\geq i$. Let $\varphi\in\PSH(X,\theta)$ be such that $\varphi\preceq_P\varphi_i$ for all $i\in I$. Then the following are equivalent:
    \begin{enumerate}
        \item\label{itm:cor:decnetdS:1} We have
            \[
                \varphi_i\xrightarrow{d_S} \varphi;
            \]
        \item\label{itm:cor:decnetdS:2} for each $k=0,\ldots,n$, we have
            \begin{equation}\label{eq:mixedmasslim}
                \lim_{j\in I}\int_X \theta_{\varphi_j}^k\wedge \theta_{V_{\theta}}^{n-k} =  \int_X \theta_{\varphi}^k\wedge \theta_{V_{\theta}}^{n-k}.
            \end{equation}
    \end{enumerate}
    If we assume furthermore that $\int_X \theta_{\varphi}^n>0$, then the above conditions are equivalent to the following:
    \begin{enumerate}[resume]
        \item\label{itm:cor:decnetdS:3} We have
            \[
                \lim_{j\in I}\int_X \theta_{\varphi_j}^n=\int_X \theta_{\varphi}^n.
            \]
    \end{enumerate}
    When $\int_X \theta_{\varphi}^n>0$ and these equivalent conditions are met, we also have
    \begin{equation}\label{eq:Pcontdecseq}
        P_{\theta}[\varphi]=\inf_{j\in I}P_{\theta}[\varphi_j].
    \end{equation}
\end{corollary}
\begin{proof}
    \ref{itm:cor:decnetdS:1} $\iff$ \ref{itm:cor:decnetdS:2}. This follows immediately from \cref{lma:varphileqpsi_metric}.

    Assume that $\int_X \theta_{\varphi}^n>0$.

    \ref{itm:cor:decnetdS:2} $\implies$ \ref{itm:cor:decnetdS:3}. This is the case $k=n$ in \eqref{eq:mixedmasslim}.

    \ref{itm:cor:decnetdS:3} $\implies$ \ref{itm:cor:decnetdS:2}.
    Assume (3). If for some $i\in I$, $\int_X \theta_{\varphi_i}^n= \int_X \theta_{\varphi}^n$, then for all $j\geq i$ we have
    \[
        \varphi\preceq_P\varphi_j\preceq_P\varphi_i.
    \]
    Hence $\int_X\theta_{\varphi_j}^n=\int_X\theta_{\varphi}^n$ by \cref{prop:mono2}. It follows from \cref{prop:Pequivchar2} that $\varphi_j\sim_P\varphi$, and (2) follows from \cref{prop:Ppresmass}. So in the sequel, we shall assume that for all $i\in I$, we have $\int_X \theta_{\varphi_i}^n>\int_X \theta_{\varphi}^n$.
    In particular, each $\varphi_i$ has positive mass.

    Let $(b_j)_{j\in I}$ be a net converging to $\infty$ such that
    \[
    b_j\in \left( 1, \left(\frac{\int_X \theta_{\varphi_j}^n}{\int_X \theta_{\varphi_j}^n-\int_X \theta_{\varphi}^n}\right)^{1/n} \right).
    \]
    This can be achieved for example by taking $b_j$ as the midpoint of this interval.
    By \cref{thm:pathoenvelope}, for each $j\in I$, we can find $\eta_j\in \PSH(X,\theta)$ such that
    \[
    b_j^{-1}\eta_j +\left(1-b_j^{-1}\right)\varphi_j \leq \varphi.
    \]
    It follows from the monotonicity theorem \cref{thm:mono} that for any $k=0,\ldots,n$,
    \[
    \int_X \theta_{\varphi}^k \wedge \theta_{V_{\theta}}^{n-k}\geq \left(1-b_j^{-1}\right)^k\int_X \theta_{\varphi_j}^k \wedge \theta_{V_{\theta}}^{n-k}.
    \]
    Taking the limit, we conclude the $\leq$ direction in \eqref{eq:mixedmasslim}. The $\geq$ direction follows from \cref{prop:mono2}.

    Finally, we argue \eqref{eq:Pcontdecseq}.
    Let $\psi_j=P_{\theta}[\varphi_j]$. It follows from \cref{cor:Psendspotentialtomodel} that $\psi_j$ is a model potential. Since the net $(\varphi_j)_{j\in I}$ is decreasing with respect to the $P$-preorder, $(\psi_j)_{j\in I}$ is decreasing pointwise. Let
    \[
        \psi=\inf_{j\in I}\psi_j\geq P_{\theta}[\varphi].
    \]
    By \cref{prop:decseqmodel}, $\psi$ is a model potential. Since $P_{\theta}[\varphi]\leq \psi\leq \psi_j$ for all $j\in I$, the monotonicity theorem \cref{thm:mono} and \cref{prop:Ppresmass} give
    \[
        \int_X \theta_{\varphi}^n=\int_X\theta_{P_{\theta}[\varphi]}^n\leq \int_X \theta_{\psi}^n\leq \lim_{j\in I}\int_X \theta_{\psi_j}^n=\lim_{j\in I}\int_X \theta_{\varphi_j}^n=\int_X \theta_{\varphi}^n.
    \]
    Hence $\psi$ has the same mass as $P_{\theta}[\varphi]$. As $\psi$ is model and $\psi\geq P_{\theta}[\varphi]$, \cref{thm:Pvarphidiffdef} gives $\psi=P_{\theta}[\varphi]$.
\end{proof}




Having understood the increasing and decreasing cases, we shall handle general convergent sequences. In fact, since $d_S$ is a pseudometric, the topology is completely determined by convergent sequences, so we do not need to consider nets in general.

\begin{proposition}\label{prop:incanddec}
    Let $\varphi_j,\varphi\in \PSH(X,\theta)$ ($j\geq 1$), $\varphi_j\xrightarrow{d_S}\varphi$. Assume that there is $\delta>0$ such that
    \[
    \int_X \theta_{\varphi_j}^n\geq \delta
    \]
    for all $j$ and the $\varphi_j$'s are all model potentials. Then up to replacing $(\varphi_j)_j$ by a subsequence, there is a decreasing sequence $(\psi_j)_j$ and an increasing sequence $(\eta_j)_j$ in $\PSH(X,\theta)$ such that
    \begin{enumerate}
        \item\label{itm:prop:incanddec:1} $\psi_j\xrightarrow{d_S}\varphi$, $\eta_j\xrightarrow{d_S}\varphi$;
        \item\label{itm:prop:incanddec:2} $\psi_j\geq \varphi_j\geq \eta_j$ for all $j$.
    \end{enumerate}
    In fact, for any $j\geq 1$, we will take
    \[
    \eta_j=\inf_{k\in \mathbb{N}}\varphi_j\land \varphi_{j+1}\land \cdots \land \varphi_{j+k},\quad
    \psi_j=\uscsup[k\geq j]\varphi_k.
    \]
\end{proposition}
\begin{proof}
    First note that $\int_X \theta_{\varphi}^n\geq \delta$ as well, due to the $d_S$-continuity of non-pluripolar masses as proved in \cref{cor:massdiffdS}. Replacing $\varphi$ by $P_{\theta}[\varphi]$, we may then assume that $\varphi$ is a model potential in $\PSH(X,\theta)_{>0}$.

    We are free to replace $(\varphi_j)_j$ by a subsequence. So we may assume that
    \begin{equation}\label{eq:conditiononvarphijtemp1}
    d_S(\varphi_j,\varphi_{j+1})\leq \frac{\delta 2^{-j-4}}{D_n C_n^{j+1}},\quad d_S(\varphi,\varphi_j)\leq \frac{\delta 2^{-j-4}}{D_n},
    \end{equation}
    where $C_n$ is the constant in \cref{cor:dsmetricdoubleineq}, $D_n$ is the constant in \cref{cor:massdiffdS}.

    In particular, by \cref{cor:massdiffdS},
    \begin{equation}\label{eq:varphijvarphimassdiffbdd}
        \left|\int_X \theta_{\varphi_j}^n-\int_X \theta_{\varphi}^n\right|\leq \delta 2^{-j-4}.
    \end{equation}

    \textbf{Step~1}. We handle the $\psi_j$'s. For each $j\geq 1$ and $k\geq 1$, by \cref{lma:dSalmostriang} we have
    \[
    \begin{aligned}
    d_S(\varphi_j,\varphi_j\lor\varphi_{j+1}\lor\cdots\lor \varphi_{j+k})
    &\leq C_n\,d_S(\varphi_j,\varphi_{j+1}\lor\cdots\lor \varphi_{j+k})\\
    &\leq C_n\,d_S(\varphi_j,\varphi_{j+1})\\
    &\quad + C_n\,d_S(\varphi_{j+1},\varphi_{j+1}\lor\cdots\lor \varphi_{j+k}).
    \end{aligned}
    \]
    By iteration, we find
    \[
    \begin{split}
        d_S(\varphi_j,\varphi_j\lor\varphi_{j+1}\lor\cdots\lor \varphi_{j+k})\leq \sum_{a=j}^{j+k-1}C_n^{a+1-j}d_S(\varphi_a,\varphi_{a+1})\\
        \leq \frac{C_n^{-j}}{D_n}\sum_{a=j}^{j+k-1}\delta 2^{-a-4}\leq \frac{\delta 2^{-j-3}}{D_n}.
    \end{split}
    \]
    Using \cref{cor:incnetdSconv}, we have
    \[
        \varphi_j\lor\varphi_{j+1}\lor\cdots\lor \varphi_{j+k}\xrightarrow{d_S}\psi_j
    \]
    as $k\to\infty$. Hence
    \begin{equation}\label{eq:dsvarphijpsijesttemp1}
        d_S(\varphi_j,\psi_j)\leq \frac{\delta 2^{-j-3}}{D_n}.
    \end{equation}
    We conclude that $\psi_j\xrightarrow{d_S}\varphi$.

    Moreover, we observe that
    \begin{equation}\label{eq:varphiexpressiontemp1}
        \varphi=\inf_{j\geq 1} P_{\theta}[\psi_j]
    \end{equation}
    Indeed, set $\tilde{\psi}_j=P_{\theta}[\psi_j]$. Then $(\tilde{\psi}_j)_j$ is a decreasing sequence of model potentials, and $d_S(\tilde{\psi}_j,\psi_j)=0$ by \cref{prop:ds0char}. Hence $\tilde{\psi}_j\xrightarrow{d_S}\varphi$. If $\chi=\inf_j\tilde{\psi}_j$, then \cref{prop:vol_limit_model} and \cref{cor:decnetdS} give $\tilde{\psi}_j\xrightarrow{d_S}\chi$. Thus $d_S(\chi,\varphi)=0$. Since both $\chi$ and $\varphi$ are model potentials, \cref{prop:ds0char} gives $P_{\theta}[\chi]=P_{\theta}[\varphi]$, hence $\chi=\varphi$, proving \eqref{eq:varphiexpressiontemp1}.

    \textbf{Step~2}. We consider the $\eta_j$'s.

    For each $j\geq 1$ and $k\geq 0$, we let
    \[
        \eta_j^k\coloneqq \varphi_j\land \cdots \land \varphi_{j+k}.
    \]
    Using \eqref{eq:conditiononvarphijtemp1}, \eqref{eq:dsvarphijpsijesttemp1} and \cref{cor:massdiffdS}, we have
    \[
        \left|\int_X\theta_{\psi_j}^n-\int_X \theta_{\varphi}^n \right|\leq \delta 2^{-j-2}
    \]
    for all $j\geq 1$. Taking \eqref{eq:varphijvarphimassdiffbdd} into account, we have
    \begin{equation}\label{eq:varphijpsijmassdiffbddtemp1}
     \left|\int_X\theta_{\varphi_j}^n-\int_X \theta_{\psi_{j-1}}^n \right|\leq \delta 2^{-j}
    \end{equation}
    for $j>1$.

    \textbf{Step~2.1}.
     We claim that for a fixed $j\geq 1$, for any $k\in \mathbb{N}$, we have
    $\eta_j^k\in \PSH(X,\theta)$ and
    \begin{equation}\label{eq:etajkmassbddbelowtemp1}
        \int_X \theta_{\eta_j^k}^n\geq \int_X \theta_{\varphi_j}^n-\sum_{a=1}^{k} \delta 2^{-j-a}.
    \end{equation}

    We argue by induction on $k\geq 0$.
    The case $k=0$ is trivial. When $k>0$, assume that the case $k-1$ is known. Then
    \[
        \begin{aligned}
            \int_X \theta_{\eta_j^{k-1}}^n+\int_X \theta_{\varphi_{j+k}}^n
            \geq \int_X \theta_{\varphi_j}^n-\sum_{a=1}^{k-1}\delta 2^{-j-a}+\int_X \theta_{\psi_{j+k-1}}^n-\delta 2^{-j-k}\\
            \geq \int_X \theta_{\varphi_j}^n-\delta 2^{-j}+\int_X \theta_{\psi_{j+k-1}}^n>\int_X \theta_{\psi_{j+k-1}}^n,
        \end{aligned}
    \]
    where the first inequality follows from the inductive hypothesis and \eqref{eq:varphijpsijmassdiffbddtemp1}.

    Observe that
    \[
    \eta_j^{k-1}\lor \varphi_{j+k}\leq \psi_{j+k-1},
    \]
    it follows from \cref{prop:landfinitecond1} that $\eta_j^k\in \PSH(X,\theta)$. By \cref{thm:diamond}, we deduce that
    \[
    \begin{aligned}
        \int_X \theta_{\eta_j^k}^n\geq & \int_X \theta_{\varphi_{j+k}}^n+  \int_X \theta_{\eta_j^{k-1}}^n-\int_X \theta_{\psi_{j+k-1}}^n\\
        \geq & \int_X \theta_{\varphi_j}^n-\sum_{a=1}^{k} \delta 2^{-j-a},
    \end{aligned}
    \]
    where the second inequality follows from the inductive hypothesis and \eqref{eq:varphijpsijmassdiffbddtemp1}. Therefore, \eqref{eq:etajkmassbddbelowtemp1} follows.


    \textbf{Step~2.2}.
    It follows by induction from \cref{prop:landpresmodel} that for any $j\geq 1$, $k\geq 0$,
    \[
        P_{\theta}\left[\eta_j^k\right]=\eta_j^k.
    \]
    By \cref{prop:vol_limit_model}, we have
    \[
        \lim_{k\to\infty} \int_X \theta_{\eta_j^k}^n=\int_X \theta_{\eta_j}^n
    \]
    for any $j\geq 1$.
    Letting $k\to\infty$ in \eqref{eq:etajkmassbddbelowtemp1}, we find that
    \[
        \int_X \theta_{\eta_j}^n\geq \int_X \theta_{\varphi_j}^n-\delta 2^{-j}>0
    \]
    for $j\geq 1$.
    Observe that we also have
    \[
        \int_X \theta_{\eta_j}^n\leq \int_X \theta_{\varphi_j}^n \leq \int_X \theta_{\psi_j}^n
    \]
    for $j\geq 1$ by \cref{thm:mono}.
    It follows from these inequalities and \eqref{eq:varphijvarphimassdiffbdd} that
    \[
        \lim_{j\to\infty}\int_X\theta_{\eta_j}^n=\int_X\theta_{\varphi}^n.
    \]
    For $k\geq 1$, we have
    \[
        \eta_j^k=\varphi_j\land\eta_{j+1}^{k-1}\leq \eta_{j+1}^{k-1}.
    \]
    Letting $k\to\infty$ gives $\eta_j\leq\eta_{j+1}$. Hence $(\eta_j)_j$ is increasing. Hence \cref{cor:incseqnppcont} gives
    \[
        \int_X \theta_{\eta}^n=\lim_{j\to\infty}\int_X \theta_{\eta_j}^n=\int_X \theta_{\varphi}^n,
    \]
    where $\eta=\uscsup[j\geq 1] \eta_j$.
    For fixed $m$, if $j\geq m$, then $\eta_j\leq\varphi_j\leq\psi_m$, whereas if $j<m$, then $\eta_j\leq\eta_m\leq\psi_m$. Thus $\eta_j$ is a candidate in the definition of $P_{\theta}[\psi_m]$ for all $j,m$, and \eqref{eq:varphiexpressiontemp1} gives $\eta\leq \varphi$.
    It follows from \cref{prop:Pequivchar2} that $\eta\sim_P \varphi$. By \cref{cor:incnetdSconv} and \cref{prop:ds0char}, we have $\eta_j\xrightarrow{d_S}\varphi$.
\end{proof}


\begin{corollary}\label{cor:completenessdS}
    Let $(\varphi_j)_{j\in I}$ be a Cauchy net (with respect to $d_S$) in $\PSH(X,\theta)$. Assume that there is $\delta>0$ such that $\int_X \theta_{\varphi_j}^n\geq \delta$ for all $j\in I$. Then $(\varphi_j)_{j\in I}$ converges with respect to $d_S$.

    In particular, if $(\varphi_j)_{j\in I}$ is a decreasing net such that $\int_X \theta_{\varphi_j}^n\geq \delta>0$ for all $j\in I$, then $(\varphi_j)_{j\in I}$ converges with respect to $d_S$.
\end{corollary}
We can obviously relax the decreasing condition to the following: the $P$-singularity types of $(\varphi_j)_{j\in I}$ are decreasing.

\begin{proof}
    In a pseudometric space, sequential completeness implies completeness for Cauchy nets. So for the first assertion, we may assume that $(\varphi_j)_j$ is a Cauchy sequence.

    We first treat the case of a decreasing sequence $(\gamma_j)_j$ with $\int_X\theta_{\gamma_j}^n\geq\delta$ for all $j$. Put $\tilde \gamma_j=P_\theta[\gamma_j]$. Then $(\tilde \gamma_j)_j$ is a decreasing sequence of model potentials, $d_S(\gamma_j,\tilde \gamma_j)=0$, and
    \[
        \int_X\theta_{\tilde \gamma_j}^n=\int_X\theta_{\gamma_j}^n\geq \delta
    \]
    by \cref{prop:Ppresmass}. By \cref{prop:vol_limit_model}, the decreasing limit $\tilde \gamma=\inf_j\tilde \gamma_j$ is model and
    \[
        \lim_{j\to\infty}\int_X\theta_{\tilde \gamma_j}^n=\int_X\theta_{\tilde \gamma}^n\geq\delta.
    \]
    Hence $\tilde \gamma_j\to\tilde \gamma$ in $d_S$ by \cref{cor:decnetdS}, and therefore $\gamma_j$ also converges in $d_S$.

    Next let $(\varphi_j)_{j>0}$ be a $d_S$-Cauchy sequence with masses bounded from below by $\delta$. Put
    \[
        \tilde{\varphi}_j=P_\theta[\varphi_j].
    \]
    By \cref{cor:Psendspotentialtomodel,ex:Pequiv,prop:ds0char}, the $\tilde{\varphi}_j$'s are model potentials, $d_S(\varphi_j,\tilde{\varphi}_j)=0$, and
    \[
        \int_X\theta_{\tilde{\varphi}_j}^n=\int_X\theta_{\varphi_j}^n\geq\delta
    \]
    by \cref{prop:Ppresmass}. Replacing $\varphi_j$ by $\tilde{\varphi}_j$, we may therefore assume that the $\varphi_j$'s are model potentials. Passing to a subsequence, we may also assume that
    \[
        d_S(\varphi_j,\varphi_{j+1})\leq C_n^{-2j},
    \]
    where $C_n$ is the constant in \cref{cor:dsmetricdoubleineq}.
    Define
    \[
    \psi_j=\uscsup[k\geq j]\varphi_k
    \]
    for each $j>0$. As in the proof of \cref{prop:incanddec} Step~1, especially \eqref{eq:dsvarphijpsijesttemp1}, we know that
    \[
    \lim_{j\to\infty} d_S(\varphi_j,\psi_j)=0.
    \]
    The sequence $(\psi_j)_j$ is decreasing and satisfies $\int_X\theta_{\psi_j}^n\geq\delta$ by \cref{thm:mono}. By the first paragraph, $(\psi_j)_j$ converges in $d_S$. Hence the chosen subsequence of $(\varphi_j)_j$ converges, and the original Cauchy sequence converges as well.

    Finally, suppose that $(\varphi_j)_{j\in I}$ is a decreasing net with masses bounded from below by $\delta$. For each $k=0,\ldots,n$, the net
    \[
        \int_X\theta_{\varphi_j}^k\wedge\theta_{V_\theta}^{n-k}
    \]
    is decreasing and bounded from below, hence Cauchy as a real-valued net. If $j'\geq j$, then $\varphi_{j'}\leq\varphi_j$, and \cref{lma:varphileqpsi_metric} expresses $d_S(\varphi_{j'},\varphi_j)$ as a finite sum of the corresponding differences of these mixed masses. For two sufficiently large indices, we choose a common upper bound in $I$ and use the triangle inequality. It follows that $(\varphi_j)_{j\in I}$ is $d_S$-Cauchy. The first assertion gives its convergence.
\end{proof}

\begin{lemma}\label{lma:dSsmallmult}
    There is a constant $C>0$ depending only on $X$ and $\theta$ such that for any $\varphi\in \PSH(X,\theta)$ satisfying that $\theta_{\varphi}$ is a Kähler current, we have
        \[
        d_{S,\theta}\Bigl((1-\epsilon)\varphi,\varphi\Bigr)\leq C \epsilon
        \]
        for small enough $\epsilon>0$ such that $(1-\epsilon)\varphi\in \PSH(X,\theta)$.
    \end{lemma}
    \begin{proof}
        After adding a constant, we may assume that $\varphi\leq 0$. Since $\theta_{\varphi}$ is a Kähler current, $(1-\epsilon)\varphi\in \PSH(X,\theta)$ for all small enough $\epsilon>0$. Fix such an $\epsilon$. Take any Kähler form $\omega$ on $X$ so that $\omega-\theta$ is still a Kähler form. Note that the choice of $\omega$ depends only on $X$ and $\theta$.

        Since $\varphi\preceq (1-\epsilon)\varphi$, by \cref{lma:varphileqpsi_metric} and the multi-linearity of the non-pluripolar product \itemref{itm:prop:npp1:3}, we can compute
        \[
        \begin{aligned}
            d_{S,\theta}\Bigl((1-\epsilon)\varphi,\varphi\Bigr)=& \frac{1}{n+1}\sum_{j=0}^n \left(\int_X \theta_{(1-\epsilon)\varphi}^j \wedge \theta_{V_{\theta}}^{n-j}-\int_X \theta_{\varphi}^j \wedge \theta_{V_{\theta}}^{n-j}\right)\\
            \leq & \frac{1}{n+1}\sum_{j=0}^n \left(\int_X \left((1-\epsilon)\theta_{\varphi}+\epsilon\omega\right)^j \wedge \theta_{V_{\theta}}^{n-j}-\int_X \theta_{\varphi}^j \wedge \theta_{V_{\theta}}^{n-j}\right)\\
            =&\sum_{m=1}^n c_m(\varphi)\epsilon^m,
        \end{aligned}
        \]
        where, for $1\leq m\leq n$,
        \[
            c_m(\varphi)=\frac{1}{n+1}\sum_{j=m}^n \binom{j}{m}
            \sum_{\ell=0}^m (-1)^{m-\ell}\binom{m}{\ell}
            \int_X \omega^\ell\wedge\theta_{\varphi}^{j-\ell}\wedge \theta_{V_{\theta}}^{n-j}.
        \]
        Note that
        \[
        \begin{aligned}
            \left|c_m(\varphi)\right| \leq & \frac{1}{n+1}\sum_{j=m}^n \binom{j}{m}
            \sum_{\ell=0}^m \binom{m}{\ell}
            \int_X \omega^\ell\wedge\theta_{\varphi}^{j-\ell}\wedge \theta_{V_{\theta}}^{n-j}\\
            \leq &  \frac{1}{n+1}\sum_{j=m}^n \binom{j}{m}
            \sum_{\ell=0}^m \binom{m}{\ell}
            \int_X \omega^\ell\wedge \theta_{V_{\theta}}^{n-\ell}
        \end{aligned}
        \]
        where we used the monotonicity theorem \cref{thm:mono}. Hence the whole expression is of the order of $\mathcal{O}(\epsilon)$, and our assertion follows.
    \end{proof}

\subsection{Convergence theorems}
\label{subsec:dS_conv}
Next we establish some important convergence theorems, allowing us to effectively manipulate the $d_S$-convergence.

\begin{lemma}\label{lma:dsconvpertV}
    Let $(\varphi_i)_{i\in I}$ be a net in $\PSH(X,\theta)$ and $\varphi\in \PSH(X,\theta)$. Assume that $\varphi_i\xrightarrow{d_S}\varphi$. Then for any $t\in (0,1]$,
    \[
    (1-t)\varphi_i+ tV_{\theta}\xrightarrow{d_S}(1-t)\varphi+ tV_{\theta}.
    \]
\end{lemma}
 When $t=1$, the sum is understood as in \cref{rmk:linearcompsh}.
\begin{proof}
    Fix $t\in (0,1]$. We write
    \[
    \varphi_{i,t}=(1-t)\varphi_i+ tV_{\theta},\quad \varphi_t=(1-t)\varphi+ tV_{\theta}
    \]
    for any $i\in I$.

    By \cref{cor:dsmetricdoubleineq}, it suffices to show that for each $j=0,\ldots,n$,
    \begin{equation}\label{eq:massconvafterpert}
    2\int_X \theta_{\varphi_{i,t}\lor \varphi_t}^j\wedge \theta_{V_{\theta}}^{n-j}-\int_X \theta_{\varphi_{i,t}}^j\wedge \theta_{V_{\theta}}^{n-j}-\int_X \theta_{\varphi_t}^j\wedge \theta_{V_{\theta}}^{n-j}\to 0.
    \end{equation}
    Observe that
    \[
    \varphi_{i,t}\lor \varphi_t=(1-t)(\varphi\lor \varphi_i)+t V_{\theta}.
    \]
    So after binomial expansion, \eqref{eq:massconvafterpert} follows from \cref{cor:dsmetricdoubleineq}.
\end{proof}

\begin{lemma}\label{lma:linearpertbyVtheta}
    Let $\varphi\in \PSH(X,\theta)$. For each $t\in (0,1)$, let $\varphi_t=(1-t)\varphi+t V_{\theta}$. Then
    \[
        \varphi_t\xrightarrow{d_S}\varphi
    \]
    as $t\to 0+$.
\end{lemma}
\begin{proof}
    Since $\varphi\preceq\varphi_t$, by \cref{lma:varphileqpsi_metric}, we need to show that for each $j=1,\dots,n$, we have
    \[
        \lim_{t\to 0+} \int_X \theta_{\varphi_t}^j\wedge \theta_{V_{\theta}}^{n-j}=\int_X \theta_{\varphi}^j\wedge \theta_{V_{\theta}}^{n-j}.
    \]
    For this purpose, we compute
    \[
        \begin{aligned}
            & \int_X \theta_{\varphi_t}^j\wedge \theta_{V_{\theta}}^{n-j}- \int_X \theta_{\varphi}^j\wedge \theta_{V_{\theta}}^{n-j}\\
            =&\left((1-t)^j-1\right)\int_X \theta_{\varphi}^j\wedge \theta_{V_{\theta}}^{n-j}\\
            &+\sum_{i=0}^{j-1} \binom{j}{i}(1-t)^i t^{j-i}\int_X\theta_{\varphi}^i\wedge \theta_{V_{\theta}}^{n-i}.
        \end{aligned}
    \]
    As $t\to 0+$, the right-hand side clearly tends to $0$.
\end{proof}


The following convergence theorem lies at the heart of the whole theory.
\begin{theorem}\label{thm:convdS}
    Let $\theta_1,\ldots,\theta_n$ be smooth closed real $(1,1)$-forms on $X$ representing big cohomology classes. Suppose that $(\varphi_j^k)_{k\in I}$ are nets in $\PSH(X,\theta_j)$ and $\varphi_j\in \PSH(X,\theta_j)$ for $j=1,\ldots,n$. We assume that $\varphi_j^k\xrightarrow{d_S}\varphi_j$ for each $j=1,\ldots,n$. Then
    \begin{equation}\label{eq:convmixedmassds}
        \lim_{k\in I}\int_X\theta_{1,\varphi_1^k}\wedge\cdots\wedge \theta_{n,\varphi_n^k}=\int_X\theta_{1,\varphi_1}\wedge\cdots\wedge \theta_{n,\varphi_n}.
    \end{equation}
\end{theorem}
\begin{proof}
Since $d_S$ is a pseudometric, in order to establish the continuity of mixed masses, it suffices to consider sequences instead of nets. So we may assume that $I=\mathbb{Z}_{>0}$ as ordered sets.

\textbf{Step~1}. We reduce to the case where $\varphi_j^k,\varphi_j$ all have positive masses and there is a constant $\delta>0$ such that for all $j=1,\ldots,n$ and $k\in I$,
	\[
		\int_X \theta_{j,\varphi_j^k}^n>\delta.
	\]
    Assuming in this case, the theorem has been proved, let us handle the general case.
    Take $t\in (0,1)$.
    By \cref{lma:dsconvpertV}, we have
    \[
    (1-t)\varphi_j^k+tV_{\theta_j}\xrightarrow{d_S}(1-t)\varphi_j+tV_{\theta_j}
    \]
    as $k\to\infty$ for each $j=1,\ldots,n$.
    Note that for all $k\in I$,
    \[
        \int_X \theta^n_{j,(1-t)\varphi_j^k+tV_{\theta_j}}\geq t^n\int_X \theta_{V_{\theta_j}}^n>0.
    \]
    Assume that we have proved the special case of the theorem, we have
    \[
    \begin{split}
    &\lim_{k\to\infty}\int_X\theta_{1,(1-t)\varphi_1^k+tV_{\theta_1}}\wedge\cdots\wedge \theta_{n,(1-t)\varphi_n^k+tV_{\theta_n}}\\
    =&\int_X\theta_{1,(1-t)\varphi_1+tV_{\theta_1}}\wedge\cdots\wedge \theta_{n,(1-t)\varphi_n+tV_{\theta_n}}.
    \end{split}
    \]
    More explicitly, for fixed $k$ the left-hand side is the value at $t$ of a polynomial $Q_k(t)$ of degree at most $n$, and the right-hand side is the value of a polynomial $Q(t)$ of degree at most $n$. Choose $n+1$ distinct numbers $t_0,\ldots,t_n\in (0,1)$. Lagrange interpolation expresses $Q_k(0)-Q(0)$ as a fixed linear combination of the numbers $Q_k(t_a)-Q(t_a)$, which tend to $0$. Hence the same formula holds at $t=0$.
    From this, \eqref{eq:convmixedmassds} follows.



    \textbf{Step~2}. We now assume that the assumptions mentioned at the beginning of Step~1 are met. Note that we may assume that $\varphi_j^k$, $\varphi_j$ are model potentials for all $j=1,\ldots,n$, $k>0$ by \cref{prop:ds0char} and \cref{cor:Psendspotentialtomodel}.

    We first handle the following monotone situation. Assume that, for each fixed $i$, the sequence $(\varphi_i^k)_k$ is either increasing or decreasing. Let $\tilde{\varphi}_i$ denote its pointwise limit, namely
\[
\tilde{\varphi}_i=\inf_{k>0}\varphi_i^k
\]
in the decreasing case and
\[
\tilde{\varphi}_i=\uscsup[k>0]\varphi_i^k
\]
in the increasing case. We assume moreover that $\tilde{\varphi}_i\sim_P\varphi_i$ for every $i$. Since mixed non-pluripolar masses are invariant under replacing each entry by a $P$-equivalent potential, it suffices in this monotone situation to prove convergence to the mixed mass of the $\tilde{\varphi}_i$'s.

After reordering the indices, we may assume that there is $i_0\in \{0,\ldots,n\}$ such that $(\varphi_i^k)_k$ is decreasing for $i\leq i_0$ and increasing for $i>i_0$. Therefore, for each $k>0$, using \cref{thm:mono}, we have
\[
\int_X \theta_{1,\varphi_1^k}\wedge\cdots\wedge\theta_{n,\varphi_n^k}\geq \int_X \theta_{1,\tilde{\varphi}_1}\wedge\cdots\wedge \theta_{i_0,\tilde{\varphi}_{i_0}}\wedge \theta_{i_0+1,\varphi^k_{i_0+1}}\wedge \cdots  \wedge\theta_{n,\varphi_n^k}.
\]
Using \cref{cor:incseqnppcont}, we therefore conclude that
\[
\varliminf_{k\to\infty}\int_X \theta_{1,\varphi_1^k}\wedge\cdots\wedge\theta_{n,\varphi_n^k}\geq \int_X \theta_{1,\tilde{\varphi}_1}\wedge\cdots\wedge\theta_{n,\tilde{\varphi}_n}.
\]
It remains to prove
	\begin{equation}\label{eq:limsup}
		\varlimsup_{k\to\infty}\int_X \theta_{1,\varphi_1^k}\wedge\cdots\wedge\theta_{n,\varphi_n^k}\leq\int_X \theta_{1,\tilde{\varphi}_1}\wedge\cdots\wedge\theta_{n,\tilde{\varphi}_n}.
	\end{equation}
By \cref{thm:mono}, for each $k>0$, we have
\[
\int_X \theta_{1,\varphi_1^k}\wedge\cdots\wedge\theta_{n,\varphi_n^k}\leq \int_X \theta_{1,\varphi_1^k}\wedge \cdots \wedge \theta_{i_0,\varphi_{i_0}^k}\wedge \theta_{i_0+1,\tilde{\varphi}_{i_0+1}}\wedge \cdots\wedge \theta_{n,\tilde{\varphi}_{n}}.
\]
When proving \eqref{eq:limsup}, we may replace $\varphi_j^k$ by $\tilde{\varphi}_j$ whenever $j>i_0$, $k>0$.
Thus, we are reduced to the case where for all $i$, $(\varphi_i^k)_k$ is decreasing.

By \cref{prop:vol_limit_model}, the masses $\int_X\theta_{i,\varphi_i^k}^n$ decrease to $\int_X\theta_{i,\tilde{\varphi}_i}^n$. Hence, thanks to \cref{lma:pathoenvelope}, for each $i=1,\ldots,n$, we may take an increasing sequence $(b^k_i)_k$ tending to $\infty$ satisfying
\[
b_i^k\in \left(1, \left( \frac{\int_X \theta_{i,\varphi_i^k}^n}{\int_X \theta_{i,\varphi_i^k}^n-\int_X \theta_{i,\tilde{\varphi}_i}^n} \right)^{1/n} \right)
\]
where the upper endpoint is understood as $\infty$ if the denominator is $0$. For each such $b_i^k$, \cref{lma:pathoenvelope} gives $\rho_i^k\in \PSH(X,\theta_i)$ with
\[
	(b^k_i)^{-1}\rho_i^k+\left(1-(b^k_i)^{-1}\right)\varphi_i^k\leq \tilde{\varphi}_i.
\]

Then by \cref{thm:mono} again,
\[
	\prod_{i=1}^n\left(1-(b^k_i)^{-1}\right)\int_X \theta_{1,\varphi_1^k}\wedge\cdots\wedge\theta_{n,\varphi_n^k}\leq \int_X \theta_{1,\tilde{\varphi}_1}\wedge\cdots\wedge\theta_{n,\tilde{\varphi}_n}.
\]
Letting $k\to\infty$, we conclude \eqref{eq:limsup}. This proves the monotone situation.

We now return to the general case. It suffices to prove that any subsequence of
\[
\int_X \theta_{1,\varphi_1^k}\wedge\cdots\wedge\theta_{n,\varphi_n^k}
\]
has a further subsequence converging to the desired limit. Starting with an arbitrary subsequence, we pass to a further subsequence, still indexed by $k$, so that the construction in \cref{prop:incanddec} applies simultaneously to the finitely many sequences $(\varphi_i^k)_k$, $i=1,\ldots,n$.

Thus, after discarding finitely many terms and relabeling, for each $i$ there are model potentials $\eta_i^k,\psi_i^k\in \PSH(X,\theta_i)$ such that $(\eta_i^k)_k$ is increasing, $(\psi_i^k)_k$ is decreasing, and
\[
\eta_i^k\leq \varphi_i^k\leq \psi_i^k,\quad
\eta_i^k\xrightarrow{d_S}\varphi_i,\quad
\psi_i^k\xrightarrow{d_S}\varphi_i.
\]
For the upper sequence, we take $\psi_i^k=P_{\theta_i}[\uscsup[\ell\geq k]\varphi_i^\ell]$; as in \eqref{eq:varphiexpressiontemp1}, this gives $\inf_k\psi_i^k=\varphi_i$. For the lower sequence, the construction in \cref{prop:incanddec} together with \cref{prop:decseqmodel} shows that the $\eta_i^k$'s are model after discarding finitely many terms, and the proof of \cref{prop:incanddec} shows that their increasing limit is $P$-equivalent to $\varphi_i$.

By \cref{thm:mono}, for every $k$,
\[
\int_X\theta_{1,\eta_1^k}\wedge\cdots\wedge\theta_{n,\eta_n^k}
\leq
\int_X\theta_{1,\varphi_1^k}\wedge\cdots\wedge\theta_{n,\varphi_n^k}
\leq
\int_X\theta_{1,\psi_1^k}\wedge\cdots\wedge\theta_{n,\psi_n^k}.
\]
The monotone situation just proved applies to both the lower and upper sequences, and both limiting mixed masses are equal to
\[
\int_X\theta_{1,\varphi_1}\wedge\cdots\wedge\theta_{n,\varphi_n}.
\]
The squeeze theorem therefore gives the desired convergence along this further subsequence. This concludes the proof.
\end{proof}


\begin{corollary}\label{cor:dsconvcrit}
    Suppose that $(\varphi_i)_{i\in I}$ is a net in $\PSH(X,\theta)$ and $\varphi\in \PSH(X,\theta)$. Then the following are equivalent:
    \begin{enumerate}
        \item\label{itm:cor:dsconvcrit:1} $\varphi_i\xrightarrow{d_S}\varphi$;
        \item\label{itm:cor:dsconvcrit:2} $\varphi_i\lor \varphi\xrightarrow{d_S}\varphi$ and
        \begin{equation}\label{eq:massconv_varphii}
        \lim_{i\in I}\int_X \theta_{\varphi_i}^j\wedge \theta_{V_{\theta}}^{n-j}=\int_X \theta_{\varphi}^j\wedge \theta_{V_{\theta}}^{n-j}
        \end{equation}
        for each $j=0,\ldots,n$;
        \item\label{itm:cor:dsconvcrit:3} for each $j=0,\ldots,n$, \eqref{eq:massconv_varphii} holds and
        \begin{equation}\label{eq:massconv_varphii2}
        \lim_{i\in I}\int_X \theta_{\varphi_i\lor \varphi}^j\wedge \theta_{V_{\theta}}^{n-j}=\int_X \theta_{\varphi}^j\wedge \theta_{V_{\theta}}^{n-j}.
        \end{equation}
    \end{enumerate}
\end{corollary}
The corollary allows us to reduce a number of convergence problems related to $d_S$ to the case $\varphi_i\geq \varphi$.
This is the most handy way of establishing $d_S$-convergence in practice.
\begin{proof}
    The equivalence between \ref{itm:cor:dsconvcrit:2} and \ref{itm:cor:dsconvcrit:3} follows directly from \cref{lma:varphileqpsi_metric}.

    \ref{itm:cor:dsconvcrit:1} $\implies$ \ref{itm:cor:dsconvcrit:2}. That $\varphi_i\lor \varphi\xrightarrow{d_S}\varphi$ follows from \cref{cor:dsmetricdoubleineq}. While \eqref{eq:massconv_varphii} follows from \cref{thm:convdS}.

    \ref{itm:cor:dsconvcrit:2} $\implies$ \ref{itm:cor:dsconvcrit:1}. By \eqref{eq:ds_biineq}, we need to show that for each $j=0,\ldots,n$, we have
    \[
    2\int_X \theta_{\varphi_i\lor \varphi}^j\wedge \theta_{V_{\theta}}^{n-j}-\int_X \theta_{\varphi}^j\wedge \theta_{V_{\theta}}^{n-j}-\int_X \theta_{\varphi_i}^j\wedge \theta_{V_{\theta}}^{n-j}\to 0.
    \]
    This follows from \cref{thm:convdS} and \eqref{eq:massconv_varphii}.
\end{proof}

\begin{corollary}\label{cor:dSconv_changetheta}
    Let $(\varphi_i)_{i\in I}$ be a net in $\PSH(X,\theta)$ and $\varphi\in \PSH(X,\theta)$. Let $\omega$ be a closed smooth positive $(1,1)$-form on $X$. Then the following are equivalent:
    \begin{enumerate}
        \item\label{itm:cor:dSconv_changetheta:1} $\varphi_i\xrightarrow{d_{S,\theta}}\varphi$;
        \item\label{itm:cor:dSconv_changetheta:2} $\varphi_i\xrightarrow{d_{S,\theta+\omega}}\varphi$.
    \end{enumerate}
\end{corollary}
In particular, there is no risk when we simply write $\varphi_i\xrightarrow{d_S}\varphi$. Similarly, consider two nets $(\varphi_i)_{i\in I}$, $(\psi_i)_{i\in I}$ in $\PSH(X,\theta)$. If both admit $d_S$-limits in $\PSH(X,\theta)$, then we can write $d_S(\varphi_i,\psi_i)\to 0$ without ambiguity.

\begin{proof}
    \ref{itm:cor:dSconv_changetheta:1} $\implies$ \ref{itm:cor:dSconv_changetheta:2}. It suffices to show that for each $j=0,\ldots,n$, we have
    \[
    \begin{split}
        2\int_X (\theta+\omega)_{\varphi_i\lor \varphi}^j\wedge (\theta+\omega)_{V_{\theta+\omega}}^{n-j}-\int_X (\theta+\omega)_{\varphi_i}^j\wedge (\theta+\omega)_{V_{\theta+\omega}}^{n-j}\\
        -\int_X (\theta+\omega)_{\varphi}^j\wedge (\theta+\omega)_{V_{\theta+\omega}}^{n-j}\to 0.
    \end{split}
    \]
    Note that this quantity is a linear combination of terms of the following form:
    \[
    \begin{split}
        2\int_X \theta_{\varphi_i\lor \varphi}^r\wedge \omega^{j-r}\wedge (\theta+\omega)_{V_{\theta+\omega}}^{n-j}-\int_X \theta_{\varphi_i}^r\wedge \omega^{j-r}\wedge (\theta+\omega)_{V_{\theta+\omega}}^{n-j}\\
        -\int_X \theta_{\varphi}^r\wedge \omega^{j-r}\wedge (\theta+\omega)_{V_{\theta+\omega}}^{n-j},
    \end{split}
    \]
    where $r=0,\ldots,j$. By \cref{thm:convdS}, it suffices to show that $\varphi \lor \varphi_i \xrightarrow{d_S}\varphi$. But this follows from \cref{cor:dsconvcrit}.

    \ref{itm:cor:dSconv_changetheta:2} $\implies$ \ref{itm:cor:dSconv_changetheta:1}. From the direction we already proved, for each $C\geq 1$, we have
    \[
        \varphi_i\xrightarrow{d_{S,\theta+C\omega}}\varphi.
    \]
    By \cref{thm:convdS}, it follows that
    \[
        \lim_{i\in I}\int_X (\theta+C\omega)_{\varphi_i}^j\wedge \theta_{V_{\theta}}^{n-j}=\int_X (\theta+C\omega)_{\varphi}^j\wedge \theta_{V_{\theta}}^{n-j}
    \]
    for all $j=0,\ldots,n$.
    It follows that
    \begin{equation}\label{eq:varphijmass_limit}
        \lim_{i\in I}\int_X \theta_{\varphi_i}^j\wedge \theta_{V_{\theta}}^{n-j}=\int_X \theta_{\varphi}^j\wedge \theta_{V_{\theta}}^{n-j}.
    \end{equation}
    By \cref{cor:dsconvcrit}, it remains to show that $\varphi_i\lor \varphi\xrightarrow{d_{S,\theta}}\varphi$. By \cref{cor:dsconvcrit} again, we know that $\varphi_i\lor \varphi\xrightarrow{d_{S,\theta+\omega}}\varphi$. So it suffices to apply \eqref{eq:varphijmass_limit} to $\varphi_i\lor \varphi$ instead of $\varphi_i$, and we conclude by \cref{lma:varphileqpsi_metric}.
\end{proof}

\begin{corollary}\label{cor:lineardsconv}
    Let $\varphi_0,\varphi_1\in \PSH(X,\theta)$. Define $\varphi_t=t\varphi_1+(1-t)\varphi_0$ for $t\in (0,1)$.
    Then
    \[
    \varphi_t\xrightarrow{d_S}\varphi_0
    \]
    as $t\to 0+$.
\end{corollary}
\begin{proof}
First note that for each $j=0,\ldots,n$,
\[
\lim_{t\to 0+}\int_X \theta_{\varphi_t}^j\wedge \theta_{V_{\theta}}^{n-j}=\int_X \theta_{\varphi_0}^j\wedge \theta_{V_{\theta}}^{n-j}.
\]
So thanks to \cref{cor:dsconvcrit}, it remains to argue that for all $j=0,\ldots,n$,
\[
\lim_{t\to 0+}\int_X \theta_{\varphi_t\lor \varphi_0}^j\wedge \theta_{V_{\theta}}^{n-j}=\int_X \theta_{\varphi_0}^j\wedge \theta_{V_{\theta}}^{n-j}.
\]
Observe that for $t\in (0,1)$, we have
\[
\varphi_t\lor \varphi_0=t(\varphi_1\lor \varphi_0)+(1-t)\varphi_0,
\]
so the desired limit follows by the same binomial expansion.
\end{proof}

We have the following sandwich criterion:
\begin{corollary}\label{cor:dsconvupplower}
    Let $(\varphi_i)_{i\in I}, (\psi_i)_{i\in I}, (\eta_i)_{i\in I}$ be three nets in $\PSH(X,\theta)$ and $\varphi\in \PSH(X,\theta)$. Assume that
    \begin{enumerate}
        \item\label{itm:cor:dsconvupplower:1} $\psi_i\preceq_P \varphi_i\preceq_P \eta_i$ for each $i\in I$;
        \item\label{itm:cor:dsconvupplower:2} $\eta_i\xrightarrow{d_S} \varphi$, $\psi_i\xrightarrow{d_S} \varphi$.
    \end{enumerate}
    Then $\varphi_i\xrightarrow{d_S}\varphi$.
\end{corollary}
\begin{proof}
    By \cref{cor:dSconv_changetheta}, we may replace $\theta$ by $\theta+\omega$, where $\omega$ is a Kähler form on $X$. In particular, we may assume that $\varphi_i,\psi_i,\eta_i\in \PSH(X,\theta)_{>0}$ for all $i\in I$.
    By \cref{prop:ds0char}, we may assume that $\varphi_i,\psi_i,\eta_i$ are model potentials for all $i\in I$ and hence $\psi_i \leq \varphi_i\leq \eta_i$ for all $i\in I$.

    It follows from \cref{thm:mono} that for each $k=0,\ldots,n$, we have
    \[
    \int_X \theta_{\psi_i}^k\wedge \theta_{V_{\theta}}^{n-k}\leq \int_X \theta_{\varphi_i}^k\wedge \theta_{V_{\theta}}^{n-k}\leq \int_X \theta_{\eta_i}^k\wedge \theta_{V_{\theta}}^{n-k}
    \]
    for all $i\in I$. By \cref{thm:convdS}, the limits, with respect to $i\in I$, of the two outer terms are $\int_X \theta_{\varphi}^k\wedge \theta_{V_{\theta}}^{n-k}$. It follows that
    \begin{equation}\label{eq:thetak_conv}
    \lim_{i\in I}\int_X \theta_{\varphi_i}^k\wedge \theta_{V_{\theta}}^{n-k}=\int_X \theta_{\varphi}^k\wedge \theta_{V_{\theta}}^{n-k}.
    \end{equation}
    By \cref{cor:dsconvcrit}, it remains to prove that $\varphi_i \lor \varphi \xrightarrow{d_S}\varphi$. By \cref{cor:dsconvcrit} and \cref{prop:Ppartialsup}, up to replacing $\psi_i$ (resp. $\varphi_i$, $\eta_i$) by $\psi_i\lor \varphi$ (resp. $\varphi_i\lor\varphi$, $\eta_i\lor\varphi$), we may assume from the beginning that $\psi_i,\varphi_i,\eta_i\geq \varphi$. Now $\varphi_i\xrightarrow{d_S}\varphi$ by \eqref{eq:thetak_conv} and \cref{lma:varphileqpsi_metric}.
\end{proof}

\begin{lemma}\label{lma:dS_res}
    Let $(\varphi_i)_{i\in I}$ be a net in $\PSH(X,\theta)$, and $\varphi\in \PSH(X,\theta)$. Fix $\lambda>0$. Then the following are equivalent:
    \begin{enumerate}
        \item\label{itm:lma:dS_res:1} $\varphi_i\xrightarrow{d_{S,\theta}}\varphi$;
        \item\label{itm:lma:dS_res:2} $\lambda\varphi_i\xrightarrow{d_{S,\lambda\theta}}\lambda\varphi$.
    \end{enumerate}
\end{lemma}
\begin{proof}
    We have $V_{\lambda\theta}=\lambda V_{\theta}$. The middle expression in \eqref{eq:ds_biineq} for the pair $(\lambda\varphi_i,\lambda\varphi)$ in the class $\lambda\theta$ is therefore $\lambda^n$ times the corresponding expression for $(\varphi_i,\varphi)$ in the class $\theta$. The equivalence follows from \cref{cor:dsmetricdoubleineq}.
\end{proof}



\begin{proposition}\label{prop:lor_dS_conv}
    Let $(\varphi_i)_{i\in I}$ (resp. $(\psi_i)_{i\in I}$) be a net in $\PSH(X,\theta)$ such that $\varphi_i \xrightarrow{d_S}\varphi\in \PSH(X,\theta)$ (resp. $\psi_i \xrightarrow{d_S}\psi\in \PSH(X,\theta)$).
    Then
    \[
    \varphi_i\lor\psi_i\xrightarrow{d_S}\varphi\lor \psi.
    \]
\end{proposition}
\begin{proof}
    Since $d_S$ is a pseudometric, we may assume that both nets are actually sequences and $I=\mathbb{Z}_{>0}$. By \cref{cor:dSconv_changetheta}, we may assume that the masses $\int_X \theta_{\varphi}^n>0$, $\int_X \theta_{\psi}^n>0$.

    By \cref{prop:ds0char} and \cref{prop:Ppartialsup}, replacing all potentials by their $P$-envelopes changes neither the hypotheses nor the conclusion. Thus we may assume that the $\varphi_j$'s, the $\psi_j$'s, $\varphi$ and $\psi$ are all model. After discarding finitely many terms, the masses of the two sequences are bounded below by a positive constant. Applying \cref{prop:incanddec} to both sequences and using the sandwich criterion \cref{cor:dsconvupplower}, it is enough to treat the case where each of $(\varphi_j)_j$ and $(\psi_j)_j$ is monotone. In this reduced situation, $(\varphi_j)_j$ (resp. $(\psi_j)_j$) converges to $\varphi$ (resp. $\psi$) almost everywhere due to \cref{prop:decPorderPenvconve}, \cref{cor:incnetdSconv} and \cref{prop:incnetmodel}.

    We handle three cases separately.

\textbf{Step~1}. Assume that both sequences are increasing.

In this case, we have $\varphi_j\lor \psi_j\nearrow \varphi\lor \psi$ almost everywhere. Therefore, $\varphi_j\lor \psi_j\xrightarrow{d_S} \varphi\lor \psi$ by \cref{cor:incnetdSconv}.

\textbf{Step~2}. Assume that one sequence, say $(\varphi_j)_j$ is increasing while the other is decreasing. Then we have
\[
\varphi_j \lor \psi \leq \varphi_j \lor \psi_j \leq \varphi \lor \psi_j.
\]
Thanks to \cref{cor:dsconvupplower}, it suffices to show that both sides converge to $\varphi\lor \psi$ with respect to $d_S$. So we reduce to the case where both sequences are decreasing.

\textbf{Step~3}. Assume that both sequences are decreasing.

In this case, due to \cref{cor:decnetdS}, it suffices to show that
\begin{equation}\label{eq:masslortemp1}
\lim_{j\to\infty} \int_X \theta_{\varphi_j\lor \psi_j}^n=\int_X \theta_{\varphi\lor \psi}^n.
\end{equation}
The $\geq$ direction follows from \cref{thm:mono}; it remains to argue the $\leq$ direction.

Thanks to \cref{lma:pathoenvelope}, we may find a sequence $(\epsilon_j)_j$ in $(0,1)$ with limit $0$ and a sequence $(\eta_j)_j$ in $\PSH(X,\theta)_{>0}$ such that
\[
    (1-\epsilon_j)\varphi_j+\epsilon_j \eta_j \leq \varphi,\quad \eta_j \leq \varphi_j.
\]
It follows that for each $j\geq 1$, we have
\[
(1-\epsilon_j)(\varphi_j\lor \psi_j)+\epsilon_j \eta_j \leq \varphi \lor \psi_j.
\]
Therefore by \cref{thm:mono},
\[
(1-\epsilon_j)^n \int_X \theta_{\varphi_j\lor \psi_j}^n\leq \int_X \theta_{\varphi \lor \psi_j}^n.
\]
Letting $j\to\infty$, we find that
\[
\lim_{j\to\infty}\int_X \theta_{\varphi_j\lor \psi_j}^n\leq \lim_{j\to\infty}\int_X \theta_{\varphi\lor \psi_j}^n.
\]
Therefore, in order to prove \eqref{eq:masslortemp1}, we may assume that one of the sequences is constant; let us say $\psi_j=\psi$ for all $j$. Repeating the same argument as before and constructing $(\epsilon_j)_j$, $(\eta_j)_j$ as above, we get
\[
(1-\epsilon_j)^n \int_X \theta_{\varphi_j\lor \psi}^n\leq \int_X \theta_{\varphi \lor \psi}^n.
\]
Letting $j\to\infty$, we conclude \eqref{eq:masslortemp1}.
\end{proof}


\begin{proposition}\label{prop:dsconvpresorder}
    Let $(\varphi_i)_{i\in I}$, $(\psi_i)_{i\in I}$ be nets in $\PSH(X,\theta)$ such that $\varphi_i\xrightarrow{d_S} \varphi\in \PSH(X,\theta)$ and $\psi_i\xrightarrow{d_S}\psi\in \PSH(X,\theta)$. Assume that $\varphi_i\preceq_P \psi_i$ for all $i\in I$. Then $\varphi\preceq_P \psi$.
\end{proposition}
\begin{proof}
    It follows from \cref{prop:lor_dS_conv} that
    \[
        \varphi_i\lor\psi_i\xrightarrow{d_S}\varphi\lor \psi.
    \]
    By \cref{lma:reform_preceqP}, we have $\varphi_i\lor\psi_i\sim_P \psi_i$ for all $i\in I$. In particular, by \cref{prop:ds0char},
    \[
        \varphi_i\lor\psi_i\xrightarrow{d_S}\psi.
    \]
    By \cref{prop:ds0char} again, $\varphi\lor\psi\sim_P\psi$ and hence $\varphi\preceq_P \psi$ by \cref{lma:reform_preceqP}.
\end{proof}

\begin{theorem}\label{thm:dSadditivity}
    Let $\theta_1$, $\theta_2$ be smooth real closed $(1,1)$-forms on $X$ representing big cohomology classes. Suppose that $(\varphi_i)_{i\in I}$ (resp. $(\psi_i)_{i\in I}$) be a net in $\PSH(X,\theta_1)$ (resp. $\PSH(X,\theta_2)$) and $\varphi\in \PSH(X,\theta_1)$ (resp. $\psi\in \PSH(X,\theta_2)$).
    Consider the following three conditions:
    \begin{enumerate}
        \item\label{itm:thm:dSadditivity:1} $\varphi_i\xrightarrow{d_{S}}\varphi$;
        \item\label{itm:thm:dSadditivity:2} $\psi_i\xrightarrow{d_{S}}\psi$;
        \item\label{itm:thm:dSadditivity:3} $\varphi_i+\psi_i\xrightarrow{d_{S}}\varphi+\psi$.
    \end{enumerate}
    Then any two of these conditions imply the third.
\end{theorem}
\begin{proof}
By \cref{cor:dSconv_changetheta}, we may assume that $\theta_1,\theta_2$ are both Kähler forms. We denote them by $\omega_1,\omega_2$ instead. Let $\omega=\omega_1+\omega_2$.

\ref{itm:thm:dSadditivity:1}+\ref{itm:thm:dSadditivity:2} $\implies$ \ref{itm:thm:dSadditivity:3}.
	It suffices to show that for each $r=0,\ldots,n$,
	\[
		2\int_X \omega_{(\varphi_i+\psi_i)\lor(\varphi+\psi)}^r\wedge \omega^{n-r}-\int_X \omega_{\varphi_i+\psi_i}^r\wedge \omega^{n-r}-\int_X \omega_{\varphi+\psi}^r\wedge\omega^{n-r}\to 0.
	\]
	Observe that for each $i\in I$,
	\[
		(\varphi_i+\psi_i)\lor (\varphi+\psi)\leq (\varphi_i\lor \varphi)+(\psi_i \lor \psi).
	\]
	The expression with the actual rooftop is non-negative and is bounded above by the same expression with this larger potential in place of $(\varphi_i+\psi_i)\lor(\varphi+\psi)$. Thus, it suffices to show that
	\[
		2\int_X \omega_{(\varphi_i\lor \varphi)+(\psi_i\lor \psi)}^r\wedge \omega^{n-r}-\int_X \omega_{\varphi_i+\psi_i}^r\wedge \omega^{n-r}-\int_X \omega_{\varphi+\psi}^r\wedge\omega^{n-r}\to 0.
	\]
	The left-hand side is a linear combination of
	\[
		2\int_X \omega_{1,\varphi_i\lor\varphi}^a\wedge \omega_{2,\psi_i\lor\psi}^{r-a}\wedge \omega^{n-r}-\int_X \omega_{1,\varphi_i}^a\wedge \omega_{2,\psi_i}^{r-a} \wedge \omega^{n-r}-\int_X \omega_{1,\varphi}^a\wedge \omega_{2,\psi}^{r-a} \wedge \omega^{n-r}
	\]
	with $a=0,\ldots,r$. Observe that $\varphi_i\lor\varphi\xrightarrow{d_S}\varphi$ and $\psi_i\lor\psi\xrightarrow{d_S}\psi$ by \cref{cor:dsmetricdoubleineq}, each term tends to $0$ by \cref{thm:convdS}.


\ref{itm:thm:dSadditivity:1}+\ref{itm:thm:dSadditivity:3} $\implies$ \ref{itm:thm:dSadditivity:2}. For each $C\geq 1$, from the direction we already proved together with \cref{lma:dS_res},
    \[
    C\varphi_i+\psi_i\xrightarrow{d_{S}} C\varphi+\psi.
    \]
By \cref{thm:convdS}, for each $j=0,\ldots,n$,
\[
\begin{split}
&\lim_{i\in I} \int_X \Bigl(C\omega_1+\omega_2+\ddc(C\varphi_i+\psi_i)\Bigr)^j\wedge \omega_2^{n-j}\\
= &\int_X \Bigl(C\omega_1+\omega_2+\ddc(C\varphi+\psi)\Bigr)^j\wedge \omega_2^{n-j}.
\end{split}
\]
Both sides are polynomials in $C$. Comparing the constant terms gives
\begin{equation}\label{eq:psii_quant_conv}
\lim_{i\in I} \int_X \omega_{2,\psi_i}^j\wedge \omega_2^{n-j}=\int_X \omega_{2,\psi}^j\wedge \omega_2^{n-j}.
\end{equation}
Therefore, \ref{itm:thm:dSadditivity:2} follows if $\psi_i\geq \psi$ for each $i$ by \cref{lma:varphileqpsi_metric}.

Next we prove the general case. By the direction that we already proved, we know that $\varphi_i+\psi\xrightarrow{d_{S}} \varphi+\psi$. By \cref{prop:lor_dS_conv}, we have
\[
\varphi_i+(\psi_i\lor \psi)\xrightarrow{d_{S}} \varphi+\psi.
\]
It follows from the special case above that $\psi_i\lor \psi\xrightarrow{d_{S}}\psi$. It follows from \eqref{eq:psii_quant_conv} and \cref{cor:dsconvcrit} that \ref{itm:thm:dSadditivity:2} holds.

\ref{itm:thm:dSadditivity:2}+\ref{itm:thm:dSadditivity:3} $\implies$ \ref{itm:thm:dSadditivity:1}. This is similar.

\end{proof}

\begin{proposition}\label{prop:decPorderPenvconve}
    Let $(\varphi_i)_{i\in I}$ be a net in $\PSH(X,\theta)$, decreasing with respect to the $P$-partial order. Assume that there is $\varphi\in \PSH(X,\theta)_{>0}$ with $\varphi_i\xrightarrow{d_S}\varphi$, then
    \begin{equation}\label{eq:decPorderPenvconve}
        \inf_{i\in I} P_{\theta}[\varphi_i]=P_{\theta}[\varphi].
    \end{equation}
\end{proposition}
\begin{proof}
    First observe that by our assumption, $(P_{\theta}[\varphi_i])_{i\in I}$ is a decreasing net. By the continuity of non-pluripolar masses \cref{cor:massdiffdS}, we have
    \[
        \lim_{i\in I} \int_X \theta_{P_{\theta}[\varphi_i]}^n=\lim_{i\in I} \int_X \theta_{\varphi_i}^n= \int_X \theta_{\varphi}^n.
    \]
    Replacing $I$ by a tail, we may therefore assume that $\varphi_i\in \PSH(X,\theta)_{>0}$ for each $i\in I$.
    Since $\varphi_i\sim_P P_\theta[\varphi_i]$, replacing $\varphi_i$ by $P_\theta[\varphi_i]$ does not change its $d_S$-class by \cref{prop:ds0char}. We may assume that each $\varphi_i$ is model and $(\varphi_i)_{i\in I}$ is decreasing. In this case, by \cref{prop:vol_limit_model}, the potential
    \[
        \psi\coloneqq \inf_{i\in I} \varphi_i
    \]
    is model in $\PSH(X,\theta)_{>0}$ with
    \[
        \int_X \theta_{\psi}^n=\lim_{i\in I} \int_X \theta_{\varphi_i}^n= \int_X \theta_{\varphi}^n.
    \]
    Therefore, by \cref{cor:decnetdS}, we find $\varphi_i\xrightarrow{d_S}\psi$. In particular, $\varphi\sim_P \psi$ by \cref{prop:ds0char}. Therefore,
    \[
        \psi=P_{\theta}[\psi]=P_{\theta}[\varphi].
    \]
    Therefore, \eqref{eq:decPorderPenvconve} follows.
\end{proof}




\begin{theorem}\label{thm:contPI}
    The map
    \[
        P_{\theta}[\bullet]_{\mathcal{I}}\colon \PSH(X,\theta)_{>0}\rightarrow \PSH(X,\theta)_{>0}
    \]
    is continuous with respect to $d_S$.
   \end{theorem}
   \begin{proof}
       Let $(\varphi_i)_{i\in \mathbb{Z}_{>0}}$ be a sequence in $\PSH(X,\theta)_{>0}$ such that $\varphi_i\xrightarrow{d_S}\varphi\in \PSH(X,\theta)_{>0}$. We want to show that
       \begin{equation}
           P_{\theta}[\varphi_i]_{\mathcal{I}}\xrightarrow{d_S}P_{\theta}[\varphi]_{\mathcal{I}}.
       \end{equation}
       We may assume that the $\varphi_i$'s and $\varphi$ are all model potentials by \cref{prop:ds0char}.


       By \cref{prop:incanddec} and \cref{cor:dsconvupplower}, we may assume that $(\varphi_i)_i$ is either increasing or decreasing.
       In the increasing case, we apply \cref{prop:incnetmodelPI} and \cref{cor:incnetdSconv}, while in the decreasing case, we apply \cref{prop:decnetmodelPI}, \cref{prop:vol_limit_model} and \cref{cor:decnetdS}.
    \end{proof}

We record the following result for later use.
\begin{lemma}\label{lma:Vthetaepsconv}
Fix a Kähler form $\omega$ on $X$.
As $\epsilon\to 0$ (from both sides), we have
\begin{equation}\label{eq:Vthetaeps_dS}
    V_{\theta+\epsilon\omega}\xlongrightarrow{d_S}V_{\theta}.
\end{equation}
\end{lemma}
\begin{proof}
There are two assertions to prove, as detailed in the two steps.

\textbf{Step~1}. We first handle the case where $\epsilon\to 0+$.

In this case \eqref{eq:Vthetaeps_dS} means
\[
    V_{\theta+\epsilon\omega}\xlongrightarrow{d_{S,\theta+\omega}}V_{\theta}
\]
as $\epsilon\to 0+$. So thanks to \cref{cor:decnetdS} and \cref{prop:vol_limit_model}, it suffices to prove the following:
\begin{equation}\label{eq:epsil_P_temp1}
    \inf_{\epsilon>0} P_{\theta+\omega}\left[ V_{\theta+\epsilon\omega} \right]=P_{\theta+\omega}\left[ V_{\theta} \right].
\end{equation}
First observe that
\[
V_{\theta}=\inf_{\epsilon>0} V_{\theta+\epsilon\omega}.
\]
In fact, the $\leq$ direction is trivial. As for the reverse inequality, it suffices to observe that the right-hand side lies in $\PSH(X,\theta)$.
Therefore, due to \cref{prop:vol_limit_model},
\[
\lim_{\epsilon \to 0+} \int_X \left( \theta+\epsilon\omega+\ddc V_{\theta+\epsilon\omega} \right)^n=\int_X \theta_{V_{\theta}}^n.
\]
Therefore, for all $\epsilon>0$ small enough, we can find $\eta_{\epsilon}\in \PSH(X,\theta+\epsilon\omega)$ and $a_{\epsilon}>0$ decreasing to $0$ so that
\[
(1-a_\epsilon)V_{\theta+\epsilon \omega} + a_\epsilon \eta_\epsilon \leq V_{\theta}.
\]
Therefore, thanks to \cref{prop:Pconc},
\[
(1-a_\epsilon)P_{\theta+\omega}\left[ V_{\theta+\epsilon \omega} \right] + a_\epsilon P_{\theta+\omega}\left[\eta_\epsilon\right] \leq P_{\theta+\omega}\left[V_{\theta}\right].
\]
Since the envelopes $P_{\theta+\omega}[\eta_\epsilon]$ have supremum $0$, they are uniformly bounded in $L^1$ by \cref{prop:L1compa}; hence $a_\epsilon P_{\theta+\omega}[\eta_\epsilon]\to 0$ in $L^1$. Letting $\epsilon\to 0$ and using \cref{prop:pshfuncdetdense}, we conclude \eqref{eq:epsil_P_temp1}.

\textbf{Step~2}. We then handle the case where $\epsilon\to 0-$.

In this case, \eqref{eq:Vthetaeps_dS} simply means
\[
V_{\theta-\epsilon\omega}\xrightarrow{d_{S,\theta}}V_{\theta}
\]
as $\epsilon\to 0+$.
But this follows from \cref{cor:incnetdSconv} if we can prove
\begin{equation}\label{eq:Vthetamepsomega_temp1}
\uscsup[\epsilon>0] V_{\theta-\epsilon\omega}=V_{\theta},
\end{equation}
where we understand that $\epsilon$ is small enough so that $\theta-\epsilon\omega$ represents a big cohomology class. Since $\{\theta\}$ is big, we can find $\varphi\in \PSH(X,\theta)$ so that
\[
\varphi\leq 0,\quad \theta_{\varphi}\geq \delta\omega
\]
for some $\delta>0$. For a small $\epsilon>0$, we then have
\[
\theta+\ddc \left( \left(1-\frac{\epsilon}{\delta}\right)V_{\theta}+ \frac{\epsilon}{\delta}\varphi\right)\geq \epsilon\omega.
\]
Therefore,
\[
\left(1-\frac{\epsilon}{\delta}\right)V_{\theta}+ \frac{\epsilon}{\delta}\varphi\leq V_{\theta-\epsilon\omega}.
\]
Letting $\epsilon\to 0+$, we then find
\[
V_{\theta}\leq \uscsup[\epsilon>0] V_{\theta-\epsilon\omega}.
\]
The reverse inequality is trivial, hence \eqref{eq:Vthetamepsomega_temp1} is established.
\end{proof}

\subsection{Continuity of invariants}\label{subsec:continv}
In this section, we prove the continuity of a few invariants of the singularities with respect to $d_S$.



    \begin{theorem}\label{thm:Lelongcont}
        Let $(\varphi_j)_{j\in I}$ be a net in $\PSH(X,\theta)$ and $\varphi_j\xrightarrow{d_S}\varphi\in \PSH(X,\theta)$. Then for any prime divisor $E$ over $X$, we have
        \begin{equation}\label{eq:convnu}
        \lim_{j\in I}\nu(\varphi_j,E)=\nu(\varphi,E).
        \end{equation}
    \end{theorem}
\begin{proof}
    First observe that since $d_S$ is a pseudometric, it suffices to prove \eqref{eq:convnu} when $I=\mathbb{Z}_{>0}$ as partially ordered sets.

    By \cref{cor:dSconv_changetheta}, after adding a Kähler form to $\theta$, we may assume that the masses of $\varphi_j$ and of $\varphi$ are bounded from below by a positive constant.

    By \cref{thm:contPI}, we may assume that $\varphi_j$ and $\varphi$ are both $\mathcal{I}$-model and hence model. When proving \eqref{eq:convnu}, we are free to pass to subsequences.

    By \cref{prop:incanddec} and a sandwich argument, we may assume that the sequence $(\varphi_j)_j$ is either increasing or decreasing. In the increasing case, the Lelong numbers $\nu(\varphi_j,E)$ decrease to a limit at least $\nu(\varphi,E)$, while $L^1$-convergence gives the opposite inequality by upper semicontinuity of Lelong numbers. In the decreasing case, \eqref{eq:convnu} follows from \cref{prop:vol_limit_model}.
\end{proof}


\begin{theorem}\label{thm:contvolu}
    Let $(\varphi_j)_{j\in I}$ be a net in $\PSH(X,\theta)$ and $\varphi\in \PSH(X,\theta)_{>0}$. Assume that $\varphi_j\xrightarrow{d_S}\varphi$. Then
    \begin{equation}\label{eq:Ivolcont}
        \vol \theta_{\varphi_j}\to \vol \theta_{\varphi},\quad \int_X \theta_{\varphi_j}^n\to \int_X \theta_{\varphi}^n.
    \end{equation}
\end{theorem}
Recall the volume is defined in \cref{def:volqpsh}. In fact, we do not have to assume the positivity of the mass of $\varphi$. The proof of the general statement is slightly more involved. See \cref{cor:contvolu2} below.
\begin{proof}
    The latter part of \eqref{eq:Ivolcont} is just a special case of \cref{thm:convdS}. It remains to prove the former part.

    We may therefore assume that $\int_X \theta_{\varphi_j}^n>0$ for all $j\in I$ after passing to a tail. Then by \cref{thm:contPI}, we have
    \[
        P_{\theta}[\varphi_j]_{\mathcal{I}}\xrightarrow{d_S}P_{\theta}[\varphi]_{\mathcal{I}}.
    \]
    Therefore, the first part of \eqref{eq:Ivolcont} follows again from \cref{thm:convdS}.
\end{proof}

Next we show that $d_S$-convergent sequences have a sort of quasi-equisingular property (cf. \eqref{eq:quasi_eq_prop}).
\begin{theorem}\label{thm:equising_cond_general}
    Let $\varphi_j,\varphi\in \PSH(X,\theta)$ ($j\in \mathbb{Z}_{>0}$). Assume that $\varphi_j\xrightarrow{d_S}\varphi$. Then for each $\lambda'>\lambda>0$, there is $j_0>0$ so that for $j\geq j_0$,
    \begin{equation}\label{eq:quasi_equi_cond}
        \mathcal{I}(\lambda' \varphi_j)\subseteq \mathcal{I}(\lambda \varphi).
    \end{equation}
\end{theorem}
\begin{proof}
Fix $\lambda'>\lambda>0$. We want to find $j_0>0$ so that for $j\geq j_0$, \eqref{eq:quasi_equi_cond} holds.

\textbf{Step~1}. We first assume that $\varphi$ has analytic singularities.

Let $\pi\colon Y\rightarrow X$ be a log resolution of $\varphi$ and let $E_1,\ldots,E_N$ be all prime divisors in the polar locus of $\varphi$ on $Y$.
Recall that by \cref{thm:valuativemulti}, a non-zero local holomorphic function $f$ lies in the right-hand side of \eqref{eq:quasi_equi_cond} if and only if for all $i=1,\ldots,N$,
\begin{equation}\label{eq:ordEif}
\ord_{E_i}(f)>\lambda \nu(\varphi,E_i)-A_X(E_i),
\end{equation}
where $A_X$ denotes the log discrepancy.
Similarly, if $f$ lies in the left-hand side of \eqref{eq:quasi_equi_cond}, then for all $i=1,\ldots,N$,
\[
\ord_{E_i}(f)> \lambda'\nu(\varphi_j,E_i)-A_X(E_i).
\]
As Lelong numbers are continuous with respect to $d_S$ by \cref{thm:Lelongcont}, we can find $j_0>0$ so that when $j\geq j_0$, $\lambda'\nu(\varphi_j,E_i)\geq \lambda \nu(\varphi,E_i)$ for all $i$. In particular, \eqref{eq:ordEif} follows.

\textbf{Step~2}. We handle the general case.

By \cref{cor:dSconv_changetheta}, we are free to increase $\theta$ and assume that $\theta_{\varphi}$ is a Kähler current.

Take a quasi-equisingular approximation $(\psi_k)_k$ of $\varphi$ in $\PSH(X,\theta)$. The existence is guaranteed by \cref{thm:qequi}.
Take $\lambda''\in (\lambda,\lambda')$. Then by definition of quasi-equisingular approximations (cf. \eqref{eq:quasi_eq_prop}), we can find $k>0$ so that
\[
    \mathcal{I}(\lambda''\psi_k)\subseteq \mathcal{I}(\lambda \varphi).
\]
Observe that $\psi_k\geq \varphi$ for all $k>0$. By \cref{prop:lor_dS_conv}, we then find that for any $k>0$, we have $\varphi_j\lor \psi_k\xrightarrow{d_S}\varphi\lor \psi_k=\psi_k$ as $j\to\infty$.
By Step~1, we can find $j_0>0$ so that for $j\geq j_0$,
\[
    \mathcal{I}\Bigl(\lambda' (\varphi_j\lor \psi_k)\Bigr)\subseteq \mathcal{I}(\lambda''\psi_k).
\]
It follows that for $j\geq j_0$,
\[
    \mathcal{I}(\lambda' \varphi_j)\subseteq  \mathcal{I}(\lambda \varphi).
\]
This completes the proof.
\end{proof}
